% Rapport MétalTech -- M2 SIA, iaelyon -- Audit et gestion de la sécurité des SI
% Compilation : latexmk -pdf Projet_MetalTech.tex (voir README.md)
\documentclass[11pt,a4paper]{article}
\usepackage[utf8]{inputenc}
\usepackage[T1]{fontenc}
\usepackage[french]{babel}
\usepackage{lmodern,microtype,amsmath}
\usepackage[top=20mm,bottom=22mm,left=20mm,right=20mm,headheight=15pt]{geometry}
\usepackage{xcolor,graphicx,tabularx,array,booktabs,multirow,enumitem,fancyhdr,ragged2e,needspace,float,xltabular,etoolbox}
\usepackage{flafter}   % une figure n'apparaît jamais avant son point d'appel
\usepackage{placeins}  % \FloatBarrier : les figures restent dans leur question
\usepackage[font=small,labelfont=bf,justification=RaggedRight,singlelinecheck=false,skip=4pt]{caption}
\usepackage[explicit]{titlesec}
\usepackage{tikz}
\usetikzlibrary{arrows.meta,positioning,fit,calc,shapes.geometric,backgrounds,babel}
\usepackage{xurl}
\usepackage{hyperref}
\usepackage{bookmark}

% ---------- Couleurs (contraste vérifié) ----------
\definecolor{navy}{HTML}{16344B}
\definecolor{teal}{HTML}{087E83}
\definecolor{pale}{HTML}{E8F3F3}
\definecolor{graytext}{HTML}{4A5761}
\definecolor{lightgray}{HTML}{F1F3F4}
\definecolor{midgray}{HTML}{9AA5AD}
\definecolor{alert}{HTML}{B3261E}
% Palette catégorielle (validée : contraste, daltonisme)
\definecolor{cA}{HTML}{2369B3}
\definecolor{cB}{HTML}{06918A}
\definecolor{cC}{HTML}{C4621B}
\definecolor{cD}{HTML}{8C4FB5}
% Niveaux de risque (toujours accompagnés d'un libellé)
\definecolor{rFaible}{HTML}{DCEFD9}
\definecolor{rModere}{HTML}{FBF0C4}
\definecolor{rEleve}{HTML}{F8D5B3}
\definecolor{rCritique}{HTML}{F2B8B5}

\hypersetup{colorlinks=true,linkcolor=navy,urlcolor=teal,citecolor=navy,
  pdftitle={Audit et gestion de la sécurité des SI -- Étude de cas MétalTech Solutions SAS},
  pdfauthor={Saleh BA-ZIGHIFAN, Abdulrahman BA GUBAIR, Abdulaziz AL-KHULAQI, Mouhamadou Moustapha DIOP, Adel ELDARAWY},
  pdfsubject={M2 SIA -- iaelyon -- Eric Disson},pdflang={fr-FR}}
\frenchsetup{SmallCapsFigTabCaptions=false,AutoSpacePunctuation=false}
\addto\captionsfrench{\renewcommand{\tablename}{Tableau}\renewcommand{\figurename}{Figure}\renewcommand{\contentsname}{Sommaire}}

\setlength{\parindent}{0pt}\setlength{\parskip}{5pt plus 1pt}
\renewcommand{\arraystretch}{1.18}
% Tableaux en 10 pt (corps du texte en 11 pt)
\AtBeginEnvironment{xltabular}{\small}
\setlist[itemize]{leftmargin=*,topsep=2pt,itemsep=2pt,parsep=0pt}
\setlist[enumerate]{leftmargin=*,topsep=2pt,itemsep=2pt,parsep=0pt}
\widowpenalty=10000 \clubpenalty=10000 \displaywidowpenalty=10000
\raggedbottom
\renewcommand{\topfraction}{0.9}\renewcommand{\bottomfraction}{0.6}
\renewcommand{\textfraction}{0.07}\renewcommand{\floatpagefraction}{0.75}
\setlength{\LTpre}{4pt}\setlength{\LTpost}{4pt}
% Pas de préférence pour couper la page avant un titre (évite les pages à moitié vides)
\makeatletter\@secpenalty=0
% Variante de \needspace sans la pénalité de -100 qui favorise les coupures prématurées
\newcommand{\needsp}[1]{\par\penalty\z@\begingroup\setlength{\dimen@}{#1}\dimen@ii\pagegoal\advance\dimen@ii-\pagetotal
  \ifdim\dimen@>\dimen@ii\ifdim\dimen@ii>\z@\vfil\fi\penalty-\@M\fi\endgroup}
\makeatother

% ---------- En-tête et pied de page ----------
\pagestyle{fancy}\fancyhf{}
\fancyhead[L]{\small\textcolor{graytext}{M2 SIA -- Audit et gestion de la sécurité des SI}}
\fancyhead[R]{\small\textcolor{graytext}{Étude de cas MétalTech Solutions SAS}}
\fancyfoot[L]{\small\textcolor{graytext}{iaelyon -- Université Jean Moulin Lyon 3 -- 2026-2027}}
\fancyfoot[R]{\small\thepage}
\renewcommand{\headrulewidth}{0.3pt}

% ---------- Titres ----------
\titleformat{\section}{\needsp{10\baselineskip}\Large\bfseries\color{navy}}{}{0pt}{#1}
\titleformat{\subsection}{\needsp{6\baselineskip}\large\bfseries\color{teal}}{}{0pt}{#1}
\titleformat{\subsubsection}{\normalsize\bfseries\color{navy}}{}{0pt}{#1}
\titlespacing*{\section}{0pt}{6pt}{6pt}
\titlespacing*{\subsection}{0pt}{10pt}{3pt}
\titlespacing*{\subsubsection}{0pt}{8pt}{2pt}

% ---------- Commandes ----------
\newcolumntype{Y}{>{\RaggedRight\arraybackslash}X}
\newcolumntype{P}[1]{>{\RaggedRight\arraybackslash}p{#1}}
\newcolumntype{C}[1]{>{\centering\arraybackslash}p{#1}}
\newcolumntype{R}[1]{>{\raggedleft\arraybackslash}p{#1}}
\newcommand{\thd}[1]{\textbf{\textcolor{navy}{#1}}}
\newcommand{\source}[1]{\par\smallskip{\footnotesize\textcolor{graytext}{\textbf{Sources :} #1}}\par}
\newcommand{\figsource}[1]{\par\vspace{-2pt}{\RaggedRight\footnotesize\textcolor{graytext}{#1}\par}}
\newcommand{\partie}[1]{\FloatBarrier\needsp{0.28\textheight}\vspace{3mm}\phantomsection\addcontentsline{toc}{part}{#1}%
  {\centering\textcolor{teal}{\rule{\linewidth}{1.2pt}}\par\vspace{1.5mm}
   {\Large\bfseries\textcolor{navy}{#1}}\par\vspace{1.5mm}\textcolor{teal}{\rule{\linewidth}{0.6pt}}\par}\vspace{2mm}}
\newcommand{\question}[3]{\FloatBarrier\needsp{13\baselineskip}\phantomsection
  \section*{#1\quad #2}\addcontentsline{toc}{section}{#1 -- #2}
  \noindent\colorbox{lightgray}{\parbox{\dimexpr\linewidth-2\fboxsep\relax}{\small\textbf{Énoncé.} #3}}\par\smallskip}
\newcommand{\retenir}[1]{\par\smallskip\noindent\colorbox{pale}{\parbox{\dimexpr\linewidth-2\fboxsep\relax}{\textbf{À retenir.} #1}}\par}
\newcommand{\tagF}{\textsf{\footnotesize[Fait]}}
\newcommand{\tagO}{\textsf{\footnotesize[Objectif]}}
\newcommand{\tagH}{\textsf{\footnotesize[Hypothèse]}}
\newcommand{\tagP}{\textsf{\footnotesize[Proposition]}}

% Styles TikZ communs
\tikzset{
  box/.style={draw=navy,rounded corners=2pt,align=center,font=\footnotesize,inner sep=3pt},
  zone/.style={draw=midgray,rounded corners=4pt,dashed,inner sep=6pt},
  arr/.style={-{Latex[length=2.2mm]},thick,draw=graytext},
  arrT/.style={-{Latex[length=2.2mm]},thick,draw=teal},
  arrB/.style={{Latex[length=2.2mm]}-{Latex[length=2.2mm]},thick,draw=teal},
  arrX/.style={-{Latex[length=2.2mm]},thick,draw=alert,dashed},
  lbl/.style={font=\scriptsize,text=graytext,align=center},
}

\begin{document}
% =====================================================================
\begin{titlepage}
\hypersetup{pageanchor=false}
\thispagestyle{empty}
\vspace*{6mm}
{\large\textcolor{navy}{\textbf{iaelyon School of Management}}}\\[3pt]
{\normalsize Université Jean Moulin Lyon 3}\\[3pt]
{\normalsize Master 2 Systèmes d'information avancés (SIA)}\par
\vspace{24mm}
\textcolor{teal}{\rule{\linewidth}{2pt}}\\[8mm]
{\huge\bfseries\textcolor{navy}{Audit et gestion de la sécurité\\[2mm]des systèmes d'information}}\\[9mm]
{\LARGE\bfseries Étude de cas : MétalTech Solutions SAS}\\[5mm]
{\large Analyse des processus, audit du système d'information et cybersécurité}\\[7mm]
\textcolor{teal}{\rule{\linewidth}{0.8pt}}
\vspace{16mm}

{\large\textbf{Réalisé par}}\\[4mm]
\begin{tabular}{@{}l@{}}
Saleh BA-ZIGHIFAN\\[2mm]
Abdulrahman BA GUBAIR\\[2mm]
Abdulaziz AL-KHULAQI\\[2mm]
Mouhamadou Moustapha DIOP\\[2mm]
Adel ELDARAWY
\end{tabular}
\vfill
\begin{tabular}{@{}ll@{}}
\textbf{Cours :} & Audit et gestion de la sécurité des systèmes d'information\\[2mm]
\textbf{Enseignant :} & Eric Disson\\[2mm]
\textbf{Année universitaire :} & 2026-2027\\[2mm]
\textbf{Date de remise :} & 29 novembre 2026
\end{tabular}
\end{titlepage}
\hypersetup{pageanchor=true}
\setcounter{page}{2}

\tableofcontents
\vfill
{\small\textcolor{graytext}{Ce rapport répond, dans l'ordre et avec leur numérotation, aux dix questions du cas pédagogique \og MétalTech Solutions SAS \fg{} (version 1, juin 2026). Les données du cas, les objectifs, les hypothèses et les propositions du groupe sont distingués. Les références entre crochets renvoient à la liste des sources, en fin de document.}}
\clearpage

% =====================================================================
\section*{Introduction}
\addcontentsline{toc}{section}{Introduction}
MétalTech Solutions SAS est une petite et moyenne entreprise (PME) industrielle de Louhans (Saône-et-Loire). Elle fabrique des pièces métalliques de précision pour l'automobile et l'aéronautique. Sa production dépend de machines connectées et de plusieurs logiciels. Une panne informatique ou une cyberattaque peut donc arrêter l'atelier, retarder les livraisons et casser la traçabilité des pièces demandée par les clients.

La partie 1 analyse les processus. La partie 2 étudie la sécurité du système d'information (SI). Pour chaque question, nous partons des données du cas, nous expliquons le problème, puis nous proposons des actions adaptées à une entreprise de 127 salariés avec une petite équipe informatique.

\subsection*{Données clés du cas}
\begin{xltabular}{\linewidth}{@{}P{58mm}Y@{}}\toprule
\thd{Donnée} & \thd{Valeur dans le cas}\\\midrule\endfirsthead
\toprule \thd{Donnée} & \thd{Valeur dans le cas}\\\midrule\endhead
\bottomrule\endfoot
Effectif, chiffre d'affaires (CA) 2025 & 127 salariés, 18,4~M€ [1, p.~4, 11, 36]\\
Résultat net, marge brute 2025 & 1,12~M€. Marge brute annoncée : 7,36~M€, soit 40~\% du CA (annexe A) [1, p.~11]\\
Clients & 85 clients actifs : 73 en automobile, 12 en aéronautique [1, p.~22]\\
Budget SI & 462~k€ par an, dont 28~k€ pour la sécurité [1, p.~27, 32]\\
Production connectée & 12 machines à commande numérique (CN), 3 machines de mesure tridimensionnelle (CMM), 25 capteurs connectés (IoT, Internet des objets) [1, p.~20]\\
Logiciels retenus & ERP SAP Business One 10.0, MES Siemens Opcenter, PLM PTC Windchill, CRM Microsoft Dynamics 365 et Power BI [1, p.~17--23]\\
Hébergement & 70~\% des traitements sur site et 30~\% dans le cloud Azure [1, p.~28]\\
\end{xltabular}
\figsource{ERP : logiciel de gestion intégré (commandes, stocks, achats, finance). MES : logiciel qui pilote et suit la production. PLM : logiciel qui gère les plans et les données techniques des produits. CRM : logiciel de gestion de la relation client.}

\needsp{19\baselineskip}
\subsection*{Synthèse SWOT}
Le cas conseille les outils d'analyse stratégique [1, p.~42]. Cette synthèse SWOT (forces, faiblesses, opportunités, menaces) reprend seulement des éléments du cas.

\par\noindent{\footnotesize\begin{tabularx}{\linewidth}{@{}YY@{}}\toprule
\thd{Forces} & \thd{Faiblesses}\\\midrule
Certifications ISO 9001, EN 9100, IATF 16949 et accréditation NADCAP (p.~4, 16). Parc de machines CN 5 axes (p.~4). ERP, MES et PLM déjà reliés (p.~18--21). Clients dans l'automobile et l'aéronautique (p.~22) &
Logiciels en fin de support et réseau de 5 ans (p.~33). Équipe SI réduite, sans SOC (p.~30, 33). Fournisseurs uniques et MRP à améliorer (p.~9, 25). Devis calculés dans Excel (p.~7)\\\midrule
\thd{Opportunités} & \thd{Menaces}\\\midrule
Stratégie 2025--2028 : croissance de 15~\% par an, aéronautique à 30~\% du CA, usine connectée (p.~35). Certification ISO 27001 prévue (p.~35) &
Hameçonnage ciblé sur les données techniques (p.~33). Rupture de titane, source unique à J+10 (p.~9). Panne d'Internet qui touche la production (p.~29). Obligations NIS2 et RGPD (p.~27, 34)\\\bottomrule
\end{tabularx}}
\figsource{Synthèse du groupe d'après le cas [1]. Données financières exclues, car plusieurs tableaux du cas sont incohérents (annexe A).}

\subsection*{Méthode et hypothèses de travail}
Le dossier se contredit parfois d'un chapitre à l'autre. Pour chaque contradiction, nous choisissons une hypothèse et nous la signalons. L'annexe A donne la liste complète avec les pages.
\begin{itemize}
\item \textbf{Logiciels.} Nous retenons le chapitre 7 (p.~17--22) : SAP Business One 10.0, Dynamics 365 et Windchill. D'autres pages citent SAP S/4HANA, Salesforce et Siemens NX (p.~6, 7, 12, 40).
\item \textbf{Chiffre d'affaires.} Nous retenons 18,4~M€ (p.~4, 11, 36). Le budget SI vaut donc 462~k€ / 18,4~M€ = 2,51~\% du CA. La page 32 calcule 3,5~\% avec un CA de 13,2~M€ : c'est le CA qui diffère, pas le calcul.
\item \textbf{Questions de cybersécurité.} Les hypothèses de chaque énoncé sont prioritaires : Windows 10, équipe SI de 3 personnes, entité importante au sens de NIS2.
\item \textbf{Chiffres proposés.} Nos objectifs, délais et coûts ne sont pas des résultats mesurés. Les exemples chiffrés sont signalés.
\end{itemize}
Dans les tableaux, quatre étiquettes distinguent les informations : \tagF{} donnée du cas, \tagO{} objectif donné par le cas, \tagH{} hypothèse du sujet ou du groupe, \tagP{} proposition du groupe.
\source{[1], p.~3--36, 38--39 et 42.}
% =====================================================================
\partie{Partie 1 -- Analyse des processus et audit des SI}

\question{9.1.1}{Cartographie du processus de fabrication}{Cartographiez le processus de fabrication en identifiant les points critiques et les goulots d'étranglement. Cartographier les types d'information par processus.}

\subsection*{Le processus en huit étapes}
Le processus va de la commande du client à l'expédition [1, p.~5]. L'ERP transmet les ordres de fabrication (OF) au MES. Le PLM fournit les plans et les listes de composants (nomenclatures) validés. La figure~\ref{fig:fab} reprend les huit étapes du cas avec l'acteur, la durée estimée, les entrées et les sorties. Nous ajoutons une décision qualité et une boucle de traitement des non-conformités (NC) : le cas décrit ce traitement dans SAP (p.~25), sans le placer dans le flux.
\subsection*{Points critiques et goulots potentiels}
Un goulot d'étranglement est un poste qui limite la production : si le travail demandé dépasse sa capacité, les pièces s'accumulent. Une étape longue n'est donc pas forcément un goulot. Le cas donne des durées, mais pas les charges ni les attentes. Nous parlons donc de goulots \emph{potentiels}.

La somme des durées estimées donne un total indicatif de 4,31 à 11,75 jours, avec des étapes successives, des journées de 8~h et une seule opération 2+. Ce total inclut les achats et la livraison : ce n'est pas un temps de travail effectif. Le devis d'exemple annonce 8 semaines [1, p.~8]. L'écart suggère des attentes (files, achats, planning). Une cartographie de la chaîne de valeur (VSM, Value Stream Mapping) sur quelques commandes réelles permet de le vérifier.


\begin{xltabular}{\linewidth}{@{}P{30mm}YP{46mm}@{}}\toprule
\thd{Étape} & \thd{Pourquoi elle peut être critique (données du cas)} & \thd{Mesure à obtenir}\\\midrule\endfirsthead
\toprule \thd{Étape} & \thd{Pourquoi elle peut être critique (données du cas)} & \thd{Mesure à obtenir}\\\midrule\endhead
\bottomrule\endfoot
2. Approvisionnement & 1 à 3 jours. Titane à J+10, fournisseurs uniques, MRP non optimisé (p.~9, 25) & OF bloqués par manque de matière, jours de rupture\\
3. Programmation & Un programme FAO et un plan validé sont nécessaires avant l'usinage (p.~5, 21) & Temps d'attente des dossiers aux méthodes\\
4--5. Usinage (opérations 1 et 2+) & Certaines machines existent en un seul exemplaire : rectifieuse, découpe laser, électroérosion (p.~5--6). Changements de série & TRS (rendement de la machine), temps de réglage, file d'attente\\
6. Contrôle & 3 CMM pour 12 machines (p.~20). Contrôle obligatoire pour l'aéronautique & Attente avant contrôle, taux de rebut\\
7. Finition & 1 à 3 jours. Procédés spéciaux sous NADCAP (p.~5, 16) & Taille des lots, attente du four ou de la ligne de traitement\\
\end{xltabular}
Pour confirmer un goulot, on compare sur 4 à 8 semaines la charge planifiée (heures d'OF) et la capacité (heures d'ouverture moins les arrêts) de chaque poste. Le MES suit aussi les files d'attente. Aucun contrôle qualité n'est supprimé pour gagner du temps.
\begin{figure}[H]
\centering
\begin{tikzpicture}[
  pstep/.style={box,fill=pale,minimum width=52mm,minimum height=10.5mm,text width=50mm},
  io/.style={font=\scriptsize,align=left,text width=38mm,text=black},
  ioR/.style={font=\scriptsize,align=left,text width=42mm,text=black},
  gp/.style={draw=cC,fill=cC!15,circle,inner sep=1pt,font=\scriptsize\bfseries,text=cC},
]
\def\dy{1.26}
\foreach \i/\t/\a in {1/{1. Réception de la commande}/{Planificateur production · 2--4 h},
                      2/{2. Approvisionnement}/{Approvisionneur · 1--3 jours},
                      3/{3. Programmation FAO}/{Technicien méthodes · 4--8 h},
                      4/{4. Usinage CN -- opération 1}/{Opérateur CN · 2--8 h},
                      5/{5. Usinage CN -- opérations 2 et +}/{Opérateur CN · 2--8 h},
                      6/{6. Contrôle qualité}/{Technicien qualité · 30 min--2 h}}{
  \node[pstep] (s\i) at (0,{-(\i-1)*\dy}) {\textbf{\t}\\[-1pt]\a};
}
\node[diamond,draw=navy,fill=white,aspect=2.4,inner sep=1pt,font=\scriptsize,align=center] (d) at (0,{-6*\dy-0.15}) {Pièces\\conformes ?};
\node[pstep] (s7) at (0,{-7*\dy-0.45}) {\textbf{7. Finition}\\[-1pt]Technicien surface · 1--3 jours};
\node[pstep] (s8) at (0,{-8*\dy-0.45}) {\textbf{8. Expédition}\\[-1pt]Magasinier expéditeur · 1--2 jours};
\foreach \x/\y in {s1/s2,s2/s3,s3/s4,s4/s5,s5/s6,s6/d,s7/s8}{\draw[arr] (\x)--(\y);}
\draw[arr] (d)--node[right,font=\scriptsize]{oui} (s7);
% boucle NC (proposition)
\node[box,draw=cC,fill=cC!8,text width=33mm,font=\scriptsize,anchor=east] (nc) at (-3.6,{-6*\dy-0.15}) {\textbf{Traitement NC} \tagP\\isoler, analyser\\retouche validée : reprise\\sinon : rebut enregistré,\\OF relancé si besoin};
\draw[arr,draw=cC] (d.west)--node[above,font=\scriptsize,text=cC]{non} (nc.east);
\draw[arr,draw=cC,dashed] (nc.west) -- (-7.35,{-6*\dy-0.15}) -- (-7.35,{-4*\dy}) -- (s5.west);
\node[font=\scriptsize,text=cC,anchor=south] at (-4.8,{-4*\dy}) {retouche validée $\rightarrow$ reprise};
% Entrées (gauche) et sorties (droite)
\foreach \i/\e/\s in {1/{Commande client, exigences (CRM, ERP)}/{OF planifié dans l'ERP},
  2/{Besoins MRP, stock, fournisseurs}/{Matière et outils prêts, lot matière},
  3/{Plans et nomenclature validés (PLM)}/{Programme FAO simulé},
  4/{OF (MES), programme, matière}/{Pièces op.~1, données machine (OPC-UA)},
  5/{}/{Pièces usinées, temps et alarmes (MES)},
  6/{Pièces, spécifications, gamme de contrôle}/{Rapport de mesure (CMM, PolyWorks)}}{
  \node[io,anchor=east] at (-2.85,{-(\i-1)*\dy}) {\e};
  \node[ioR,anchor=west] at (2.85,{-(\i-1)*\dy}) {\s};
}
\node[io,anchor=east] at (-2.85,{-7*\dy-0.45}) {Pièces conformes, paramètres procédé};
\node[ioR,anchor=west] at (2.85,{-7*\dy-0.45}) {Pièces finies, enregistrements des procédés spéciaux (NADCAP)};
\node[io,anchor=east] at (-2.85,{-8*\dy-0.45}) {Pièces validées, commande, adresse};
\node[ioR,anchor=west] at (2.85,{-8*\dy-0.45}) {Bon de livraison, facture (ERP), suivi transport};
% En-têtes de colonnes
\node[font=\scriptsize\bfseries,text=navy] at (-4.75,0.95) {Entrées};
\node[font=\scriptsize\bfseries,text=navy] at (0,0.95) {Étape, acteur et durée estimée};
\node[font=\scriptsize\bfseries,text=navy] at (4.95,0.95) {Sorties et systèmes};
% Goulots potentiels
\foreach \n in {s2,s3,s6,s7}{\node[gp] at ($(\n.east)+(0.12,0.33)$) {G};}
\end{tikzpicture}
\caption{Processus de fabrication de MétalTech : huit étapes, entrées et sorties.}
\label{fig:fab}
\figsource{D'après [1], p.~5--6, 16--21 et 25. \textcolor{cC}{\textbf{G}} : goulot potentiel à confirmer par une mesure de charge et de capacité. La boucle NC est une proposition du groupe. FAO : fabrication assistée par ordinateur.}
\end{figure}

\subsection*{Cartographie des types d'information par processus}
L'ERP est le centre d'échange [1, p.~17]. Les flux passent par le logiciel d'intégration SAP PI/PO et par des connecteurs (p.~18). Sigles du tableau : MRP, calcul des besoins en matières. ECN, demande de modification technique. CAO et FAO, conception et fabrication assistées par ordinateur. EDI, échange électronique de documents. WMS, logiciel de gestion d'entrepôt. RFID, identification par radio.

\begin{xltabular}{\linewidth}{@{}P{24mm}YYP{37mm}@{}}\toprule
\thd{Processus} & \thd{Informations reçues} & \thd{Informations produites} & \thd{Systèmes et fréquence}\\\midrule\endfirsthead
\toprule \thd{Processus} & \thd{Informations reçues} & \thd{Informations produites} & \thd{Systèmes et fréquence}\\\midrule\endhead
\bottomrule\endfoot
Commercial & Demande client, prix, délais, capacité & Devis, commandes, données client & CRM D365, ERP, Excel. Commandes vers la production : quotidien\\
Méthodes et bureau d'études & Cahier des charges, demandes de modification (ECN) & Plans, nomenclatures, gammes, programmes FAO & PLM, CAO, ERP. Nomenclatures : quotidien\\
Achats & Besoins nets (MRP), délais fournisseurs & Commandes fournisseurs, relances & ERP, EDI (8 fournisseurs), Excel\\
Logistique et stock & Réceptions, stocks, expéditions & Disponibilités, lots, palettes & ERP, WMS, RFID (en test). Temps réel\\
Production & OF, plans, stocks, affectation des opérateurs & Avancement, temps, mesures, alarmes & MES, machines. Temps réel (moins de 5 min)\\
Qualité & Mesures, spécifications, alertes du MES & Validation, NC, traçabilité des lots & Module qualité de SAP, PolyWorks, Quality Portal. Temps réel\\
Finance & Coûts réels, factures, paiements & Marge par affaire, budget, reporting & ERP, Power BI. Hebdomadaire, comptabilité la nuit\\
RH & Compétences, présence, recrutements & Planning d'affectation des opérateurs & SAP HCM, TalentSoft. Hebdomadaire\\
\end{xltabular}
La figure~\ref{fig:flux} montre ces échanges entre les applications [1, p.~13--25].
\begin{figure}[H]
\centering
\begin{tikzpicture}[
  app/.style={box,fill=white,minimum width=27mm,minimum height=9mm,text width=26mm},
  hub/.style={box,fill=pale,minimum width=34mm,minimum height=11mm,line width=0.9pt},
  el/.style={font=\scriptsize,fill=white,inner sep=1pt,align=center},
  appW/.style={box,fill=white,minimum height=9mm,text width=44mm},
]
\node[hub] (erp) at (0,0) {\textbf{ERP SAP Business One}\\hub, via SAP PI/PO};
\node[app] (crm) at (-6.1,0) {CRM Dynamics 365\\(commercial)};
\node[app] (plm) at (6.1,0) {PLM Windchill\\(méthodes, BE)};
\node[appW] (bi) at (0,2.45) {Power BI (direction, pilotage)\\données ERP + MES (API, 5 min)};
\node[app] (mes) at (0,-2.55) {MES Opcenter\\(production, qualité)};
\node[app] (ot) at (0,-5.1) {12 CN, 3 CMM,\\25 capteurs IoT};
\node[app] (xl) at (-5.7,2.45) {Excel\\(devis, planning fourn.)};
\node[app] (edi) at (5.7,2.45) {EDI / portails\\8 fournisseurs};
\node[app] (hr) at (5.7,-2.55) {SAP HCM, TalentSoft\\(RH)};
\draw[arrB] (crm)--node[el,above=2pt,pos=0.45]{clients, commandes,\\facturation} (erp);
\draw[arrB] (plm)--node[el,above=2pt,pos=0.45]{nomenclatures,\\plans (quotidien)} (erp);
\draw[arrT] (erp)--node[el,right=2pt]{extractions} (bi);
\draw[arrB] (erp)--node[el,right=3pt,align=left]{OF, gammes / avancement,\\stocks (< 5 min)} (mes);
\draw[arrB] (mes)--node[el,right=3pt,align=left]{consignes / mesures,\\alertes (OPC-UA, MQTT)} (ot);
\draw[arrB,dashed,draw=midgray] (xl)--node[el,pos=0.5]{saisie manuelle} (erp);
\draw[arrB] (edi)--node[el,pos=0.5]{commandes\\fournisseurs} (erp);
\draw[arrT] (hr)--node[el,above=2pt]{planning d'affectation} (mes);
\end{tikzpicture}
\caption{Flux d'information entre les applications de MétalTech (existant).}
\label{fig:flux}
\figsource{D'après [1], p.~17--24. Trait plein : intégration décrite dans le cas. Pointillé gris : flux manuel. Les sens et fréquences sont ceux du cas.}
\end{figure}

\subsection*{Preuves d'audit à rechercher}
Un audit s'appuie sur des preuves fiables et traçables : documents, entretiens, observations et tests [5, diap.~28--29]. Nous proposons de collecter :
\begin{itemize}
\item quelques OF réels suivis de bout en bout, avec leurs dates dans l'ERP et le MES, pour mesurer les attentes ;
\item une comparaison des quantités ERP et MES sur un échantillon d'OF, et des versions de plans du PLM avec les programmes utilisés ;
\item un test de traçabilité d'une pièce aéronautique, du lot matière au lot client (procédure PC-QTR-005, p.~40), et les enregistrements des NC et actions correctives (PC-QNC-004, p.~40).
\end{itemize}
\source{[1], p.~5--6, 8, 13--25, 40 · [5], diap.~28--29.}
% ---------------------------------------------------------------------
\question{9.1.2}{Optimisation du flux de production}{Proposez des améliorations pour optimiser le flux de production et réduire les temps de changement de série.}

Un changement de série est le passage d'une fabrication à une autre sur la même machine : changer les outils, régler, charger le programme et valider la première pièce. Le cas ne donne pas sa durée réelle. Il faut donc la mesurer avant de fixer un objectif. Le cas vise un TRS de 75~\% et moins de 2~\% de rebut \tagO{} [1, p.~20] : ce sont des objectifs, pas des mesures.

Le TRS (taux de rendement synthétique) mesure l'efficacité d'une machine : $\text{TRS} = \text{disponibilité} \times \text{performance} \times \text{qualité}$. Le SMED (Single Minute Exchange of Die) raccourcit un changement de série en préparant à l'avance, machine en marche, tout ce qui peut l'être. Le 5S organise le poste en cinq étapes : trier, ranger, nettoyer, standardiser, maintenir.

\subsection*{Nos propositions}
\begin{xltabular}{\linewidth}{@{}P{30mm}YP{40mm}@{}}\toprule
\thd{Action} & \thd{Mise en place concrète} & \thd{Résultat attendu}\\\midrule\endfirsthead
\toprule \thd{Action} & \thd{Mise en place concrète} & \thd{Résultat attendu}\\\midrule\endhead
\bottomrule\endfoot
SMED & Filmer 3 à 5 changements. Séparer les opérations \emph{internes} (machine arrêtée) et \emph{externes} (possibles pendant que la machine tourne). Faire passer le plus possible d'opérations en externe, puis standardiser & Arrêt plus court entre deux séries\\
Préparation des outils & Préparer un chariot avant l'arrêt : outils, porte-pièces, programme vérifié, documents. Prérégler les outils hors de la machine & Moins de recherches et d'attente\\
5S et standards & Ranger et identifier les emplacements. Rédiger une fiche de réglage par famille de pièces & Moins d'erreurs et de pertes de temps\\
Planification dans le MES & Regrouper les séries compatibles (même matière, mêmes outils) sans retarder les commandes urgentes & Moins de changements inutiles\\
Maintenance préventive & Utiliser les données machines du MES pour planifier l'entretien hors des périodes chargées & Moins de pannes imprévues\\
Contrôle de la première pièce & Garder ce contrôle obligatoire, le standardiser et enregistrer le résultat dans le MES & Moins de rebuts, sans baisser la qualité\\
\end{xltabular}

\subsection*{Pilote, mesures et suivi}
\begin{itemize}
\item \textbf{Mois 1} : mesurer la situation de départ sur une machine pilote (durée et nombre de changements, TRS, rebuts). Responsables : méthodes, chef d'atelier usinage.
\item \textbf{Mois 2--3} : appliquer SMED et 5S, former les opérateurs, rédiger les fiches de réglage (chef d'atelier, opérateurs).
\item \textbf{Mois 4} : mesurer à nouveau dans les mêmes conditions et comparer (méthodes, qualité). \textbf{Mois 4--6} : étendre aux autres machines si le gain est confirmé et adapter le planning du MES (production, SI).
\end{itemize}
Indicateurs : temps moyen de changement, nombre de changements par semaine, TRS de la machine pilote, taux de rebut, taux de première pièce bonne. Le gain se calcule par rapport à la mesure de départ.

\subsection*{Exemple chiffré (illustratif)}
\textit{Hypothèses du groupe, pas des données du cas} : une machine, 8 changements par semaine de 90~min, 46 semaines par an, 80~€ par heure machine. L'arrêt actuel vaut 8 $\times$ 90~min = 12~h par semaine. Une baisse de 20~\% libère 2,4~h par semaine, soit 110,4~h par an, c'est-à-dire 110,4~h $\times$ 80~€ = 8\,832~€ de capacité théorique. Ce n'est ni une économie de trésorerie ni un bénéfice certain : le gain n'est réel que si ces heures servent à produire des pièces vendues. Le coût du pilote (formation, chariot, préréglage) reste à estimer.
\retenir{commencer petit, sur une seule machine, mesurer avant et après, puis généraliser seulement si le gain est prouvé. La qualité et les contrôles obligatoires ne sont jamais réduits.}
\source{[1], p.~5--6, 20, 30, 33. SMED et 5S : méthodes proposées par le groupe.}
% ---------------------------------------------------------------------
\question{9.1.3}{Amélioration du processus commercial}{Analysez le processus commercial et suggérez des axes pour améliorer le taux de conversion des devis.}

\subsection*{Le parcours commercial}
Le cas décrit huit étapes, de la prospection à l'analyse des revenus [1, p.~7]. La page~7 cite Salesforce, SAP SD et SAP FI, mais le chapitre~7 décrit Dynamics 365 et SAP Business One (p.~19, 22) : nous retenons ces derniers. Le DSO est le délai moyen de paiement des clients (45 jours, p.~24).
\begin{figure}[htp]
\centering
\begin{tikzpicture}[
  stage/.style={box,fill=pale,minimum width=35mm,minimum height=15mm,text width=34mm},
  pain/.style={font=\scriptsize,text=alert,align=center,text width=37mm},
]
\foreach \i/\t/\a in {1/{Prospection}/{Chef des ventes\\CRM},
                      2/{Qualification}/{Chef des ventes\\CRM},
                      3/{Devis}/{Ingénieur commercial\\ERP + Excel},
                      4/{Négociation}/{Ing. commercial + direction\\email, réunions}}{
  \node[stage] (c\i) at ({(\i-1)*4.0},0) {\textbf{\i. \t}\\[1pt]\a};
}
\foreach \i/\t/\a in {5/{Commande}/{Assistant commercial\\ERP (ventes)},
                      6/{Suivi livraison}/{Assistant commercial\\ERP + portail},
                      7/{Facturation}/{Service facturation\\ERP (finance)},
                      8/{Revenus, fidélisation}/{Direction commerciale\\CRM + Power BI}}{
  \node[stage] (c\i) at ({(\i-5)*4.0},-4.1) {\textbf{\i. \t}\\[1pt]\a};
}
\foreach \x/\y in {c1/c2,c2/c3,c3/c4,c5/c6,c6/c7,c7/c8}{\draw[arr] (\x)--(\y);}
\draw[arr] (c4.east) -- ++(0.3,0) |- (-2.15,-2.55) |- (c5.west);
\node[pain,below=1mm of c2] {risque : critères de\\qualification non décrits};
\node[pain,below=1mm of c3] {double saisie Excel et ERP,\\quantités non indiquées\\dans l'exemple};
\node[pain,below=1mm of c4] {suivi des motifs de perte\\non décrit (à vérifier)};
\node[pain,below=1mm of c5] {capacité et stock\\à vérifier avant le délai};
\node[pain,below=1mm of c7] {relances (DSO 45 j\\indiqué p.~24)};
% repère du taux de conversion
\draw[cB,line width=1.2pt,{Bar[]}-{Bar[]}] ($(c3.north west)+(0,0.25)$) -- node[above,font=\scriptsize,text=cB]{taux de conversion des devis : devis gagnés / devis clos} ($(c4.north east)+(0,0.25)$);
\end{tikzpicture}
\caption{Parcours commercial de MétalTech et points de vigilance.}
\label{fig:com}
\figsource{D'après [1], p.~7--8, 19, 22, 24. Outils : hypothèse Dynamics 365 (CRM) et SAP Business One (ERP). En rouge : faits du cas (double saisie, quantités absentes, DSO) et risques que le cas ne permet pas de confirmer, à vérifier sur place.}
\end{figure}

\subsection*{Analyse : ce que montrent les données}
\begin{itemize}
\item \textbf{Devis préparés avec SAP et Excel} (p.~7). Coûts et marges sont calculés dans Excel : doubles saisies, plusieurs versions d'un devis, risque d'erreur.
\item \textbf{Devis d'exemple ambigu} (p.~8). Le projet porte sur un lot de 500 pièces, mais les lignes n'ont pas de quantité. Le total HT (281,70~€) est la somme des trois prix unitaires affichés (125,00 + 89,50 + 67,20~€). Le montant du lot ne peut donc pas être vérifié. Un modèle de devis généré par l'ERP réduit ce risque.
\item \textbf{Le taux d'environ 35~\%} est celui du pipeline d'opportunités du CRM (p.~22), et le pipeline moyen de 1,2~M€ est le montant des opportunités ouvertes, pas un CA annuel. Le cas ne donne pas le taux de conversion des devis.
\item \textbf{Le CRM suit déjà} le délai moyen de réponse, le taux de réponse aux appels d'offres et la marge par affaire (p.~22). Le cas ne dit pas si les motifs de perte et la marge prévue de chaque devis sont suivis.
\end{itemize}

\subsection*{Définir le taux de conversion des devis}
Nous proposons deux indicateurs, calculés sur les devis \emph{clos} (gagnés ou perdus) d'une même période : \textbf{en nombre}, devis gagnés / (gagnés + perdus) $\times$ 100 ; \textbf{en valeur}, montant des devis gagnés / montant des devis clos $\times$ 100. Gagner beaucoup de petits devis et perdre les gros donne un bon taux en nombre et un mauvais taux en valeur. La première action est de calculer ces deux taux sur les 12 derniers mois, pour connaître le vrai point de départ.

\subsection*{Axes d'amélioration}
\begin{xltabular}{\linewidth}{@{}P{29mm}YP{30mm}P{28mm}@{}}\toprule
\thd{Problème ou risque} & \thd{Action proposée} & \thd{Responsable} & \thd{Indicateur}\\\midrule\endfirsthead
\toprule \thd{Problème ou risque} & \thd{Action proposée} & \thd{Responsable} & \thd{Indicateur}\\\midrule\endhead
\bottomrule\endfoot
Critères de qualification non décrits & Ajouter dans le CRM des critères de qualification : besoin, budget, délai, faisabilité technique & Chef des ventes & Part des demandes qualifiées\\
Devis lents ou ambigus & Créer un modèle de devis unique dans l'ERP, avec quantités, prix unitaires, coûts, marge, délai et validation technique. Garder Excel seulement en secours & Directrice commerciale, SI & Délai de réponse, erreurs détectées\\
Délai promis peu fiable & Vérifier le stock et la capacité (ERP, MES) avant d'annoncer un délai & ADV (administration des ventes), planification & Respect des délais promis\\
Opportunités oubliées & Programmer des relances automatiques dans Dynamics 365, par exemple à J+7 et J+21 & Ingénieurs commerciaux & Devis relancés à temps\\
Suivi des motifs de perte non décrit & Rendre obligatoire un motif de perte : prix, délai, technique, concurrent & Ingénieurs commerciaux & Part des devis perdus avec motif\\
Pilotage & Créer un tableau Power BI : taux en nombre et en valeur, par secteur et par client & Direction commerciale & Revue mensuelle\\
\end{xltabular}
Ces actions utilisent les outils en place (CRM, ERP, Power BI), sous réserve des licences et du paramétrage. Nous proposons un test de trois mois avec l'équipe commerciale (2 ingénieurs commerciaux et 4 assistantes, p.~13), puis une décision selon les résultats.

\textbf{Effet attendu (scénario illustratif, pas une prévision).} Avec le même pipeline de 1,2~M€ et un taux en valeur qui passe de 35~\% à 40~\%, le montant attendu passerait de 420~k€ à 480~k€, soit 60~k€ de commandes en plus. Ces 60~k€ sont un CA potentiel, ni garanti ni égal à un bénéfice : il faut déduire les coûts. Le but n'est pas de gagner des devis en baissant les prix.
\source{[1], p.~7--8, 13, 17--24.}
% ---------------------------------------------------------------------
\question{9.1.4}{Résilience de la supply chain}{Évaluez la supply chain et proposez des actions pour renforcer la résilience face aux ruptures d'approvisionnement.}

\subsection*{Évaluation de la situation}
La supply chain (chaîne d'approvisionnement) relie fournisseurs, réception, stock, production et clients [1, p.~9]. MétalTech annonce un double sourcing (deux fournisseurs) pour les matières stratégiques, classe ses articles par importance (méthode ABC) et utilise un Kanban électronique qui signale les besoins. Son stock de sécurité dépend du délai et de la demande maximale (p.~9). Le cas cite pourtant des faiblesses (p.~25) : calcul des besoins (MRP) à améliorer, seuil de commande fondé sur l'expérience, fournisseurs uniques, pas de visibilité en temps réel. Le planning fournisseurs est suivi dans Excel (p.~24).
\begin{figure}[H]
\centering
\begin{tikzpicture}[
  st/.style={box,fill=pale,minimum width=22mm,minimum height=11mm,text width=21mm},
  info/.style={box,fill=white,draw=cB,text width=30mm,font=\scriptsize},
]
\node[st] (f) at (0,0) {Fournisseurs\\(5 + divers)};
\node[st] (r) at (2.85,0) {Réception\\contrôle};
\node[st] (s) at (5.7,0) {Stock\\(WMS)};
\node[st] (p) at (8.55,0) {Production\\(OF)};
\node[st] (e) at (11.4,0) {Expédition};
\node[st] (c) at (14.25,0) {Clients\\(85)};
\foreach \x/\y in {f/r,r/s,s/p,p/e,e/c}{\draw[arr,line width=1.1pt] (\x)--(\y);}
\node[lbl,above=1mm of r] {quantité, qualité};
\node[lbl,above=1mm of e] {route, avion};
\node[info] (mrp) at (5.7,-2.3) {MRP de l'ERP :\\besoins nets, propositions\\de commande};
\node[info] (edi) at (1.4,-2.3) {EDI : 8 fournisseurs\\Excel : planning, reliquats};
\node[info] (kpi) at (11.4,-2.3) {Power BI : délai fournisseur,\\taux de rupture, litiges};
\draw[arrT,dashed] (mrp)--(edi);
\draw[arrT,dashed] (edi)--(f);
\draw[arrT,dashed] (s)--(mrp);
\draw[arrT,dashed] (p)--(mrp);
\draw[arrT,dashed] (p.south)--(kpi);
\end{tikzpicture}
\caption{Flux physiques (traits pleins) et flux d'information (pointillés) de la supply chain.}
\label{fig:sc}
\figsource{D'après [1], p.~9--10 et 24--25.}
\end{figure}

\subsection*{Fournisseurs et concentration des achats}
\begin{xltabular}{\linewidth}{@{}P{31mm}P{22mm}R{18mm}R{12mm}R{14mm}Y@{}}\toprule
\thd{Fournisseur} & \thd{Matière} & \thd{Achats (k€/an)} & \thd{Part} & \thd{Délai} & \thd{Sources pour cette matière dans le tableau du cas}\\\midrule\endfirsthead
\toprule \thd{Fournisseur} & \thd{Matière} & \thd{Achats (k€/an)} & \thd{Part} & \thd{Délai} & \thd{Sources pour cette matière dans le tableau du cas}\\\midrule\endhead
\bottomrule\endfoot
Aciers du Rhône & Acier carbone & 1\,850 & 28~\% & J+3 & Deux (avec Aciers du Nord)\\
Aciers du Nord & Acier carbone & 1\,320 & 20~\% & J+4 & Deux (avec Aciers du Rhône)\\
Aluminium Sud-Est & Aluminium & 1\,420 & 21~\% & J+2 & Une seule\\
Titanium Works GmbH & Titane & 980 & 15~\% & J+10 & Une seule\\
Special Metals SA & Alliages spéciaux & 680 & 10~\% & J+7 & Une seule\\
Autres & Divers & 380 & 6~\% & Variable & --\\\midrule
\textbf{Total du tableau} & & \textbf{6\,630} & \textbf{100~\%} & & \\
\end{xltabular}
\figsource{Données du cas [1, p.~9] \tagF. Le total de 6\,630~k€ est celui de ce tableau. Ce n'est pas le total des achats de l'entreprise : le compte de résultat indique 7\,200~k€ d'achats (p.~11). La dernière colonne est notre lecture du tableau \tagH. D'autres fournisseurs peuvent exister sans être cités.}

Dans ce tableau, seul l'acier a deux sources (48~\% du total, somme des parts arrondies). L'aluminium, le titane et les alliages (46~\%) n'ont qu'un fournisseur cité chacun. Le titane cumule source unique, délai le plus long (J+10) et usage aéronautique : c'est le point le plus fragile. La figure~\ref{fig:crit} place les fournisseurs selon leur délai et leur poids.
\begin{figure}[H]
\centering
\begin{tikzpicture}[x=0.78cm,y=0.19cm]
\fill[rFaible!60] (0,0) rectangle (5,12);
\fill[rModere!70] (5,0) rectangle (12,12);
\fill[rModere!70] (0,12) rectangle (5,30);
\fill[rEleve!70] (5,12) rectangle (12,30);
\draw[midgray] (0,0) rectangle (12,30);
\draw[midgray,dashed] (5,0)--(5,30); \draw[midgray,dashed] (0,12)--(12,12);
\foreach \x in {0,2,4,6,8,10,12}{\draw[graytext] (\x,0)--(\x,-0.6) node[below,font=\scriptsize]{J+\x};}
\foreach \y in {0,10,20,30}{\draw[graytext] (0,\y)--(-0.12,\y) node[left,font=\scriptsize]{\y\,\%};}
\node[font=\scriptsize,below] at (6,-3.2) {Délai de livraison du fournisseur};
\node[font=\scriptsize,rotate=90] at (-1.25,15) {Part des achats};
\node[font=\scriptsize\bfseries,text=graytext,anchor=north east] at (11.9,29.6) {À traiter en priorité};
\node[font=\scriptsize,text=graytext,anchor=south east] at (11.9,0.5) {Délai long};
\node[font=\scriptsize,text=graytext,anchor=north west] at (0.1,29.6) {Poids élevé};
\node[font=\scriptsize,text=graytext,anchor=south west] at (0.1,0.5) {Faible};
% source double = carré, source unique = disque
\foreach \x/\y/\n/\pos in {3/28/{Aciers du Rhône}/right, 4/20/{Aciers du Nord}/right}{
  \node[draw=cA,fill=cA,minimum size=6pt,inner sep=0pt] at (\x,\y) {};
  \node[font=\scriptsize,anchor=west] at (\x+0.2,\y) {\n};}
\foreach \x/\y/\n in {2/21/{Aluminium Sud-Est}, 10/15/{Titanium Works}, 7/10/{Special Metals}}{
  \node[draw=cC,fill=cC,circle,minimum size=7pt,inner sep=0pt] at (\x,\y) {};}
\node[font=\scriptsize,anchor=south] at (2,21.8) {Aluminium Sud-Est};
\node[font=\scriptsize,anchor=south] at (10,15.9) {\textbf{Titanium Works}};
\node[font=\scriptsize,anchor=north] at (7,9.1) {Special Metals};
% légende
\node[draw=cA,fill=cA,minimum size=6pt,inner sep=0pt] at (13,24) {}; \node[font=\scriptsize,anchor=west] at (13.3,24) {source double};
\node[draw=cC,fill=cC,circle,minimum size=7pt,inner sep=0pt] at (13,20) {}; \node[font=\scriptsize,anchor=west] at (13.3,20) {source unique};
\end{tikzpicture}
\caption{Matrice de criticité des fournisseurs (délai et poids dans les achats).}
\label{fig:crit}
\figsource{Positions : données du cas [1, p.~9]. Seuils (J+5 et 12\,\%) et zones : proposition du groupe. Une analyse complète ajoutera la difficulté de qualifier une autre source (exigences aéronautiques).}
\end{figure}

\subsection*{Plan d'actions}
\begin{xltabular}{\linewidth}{@{}P{20mm}YP{34mm}@{}}\toprule
\thd{Délai} & \thd{Action} & \thd{Responsable}\\\midrule\endfirsthead
\toprule \thd{Délai} & \thd{Action} & \thd{Responsable}\\\midrule\endhead
\bottomrule\endfoot
0--1 mois & Lister les références critiques et les matières à source unique. C'est aussi l'action PCA-PAP-04 du PCA [8, §9.3] & Achats, qualité\\
1--3 mois & Recalculer avec l'historique SAP le point de commande et le stock de sécurité des articles A et des matières critiques. Remplacer le point de commande empirique & Achats, logistique\\
3--6 mois & Qualifier une deuxième source pour le titane, puis pour les alliages et l'aluminium : audit qualité, échantillons, accord du client si la pièce l'exige & Achats, qualité\\
3--6 mois & Étendre l'EDI ou un portail fournisseur pour suivre les confirmations et les retards. Sortir le planning fournisseurs d'Excel & SI, achats\\
6--12 mois & Tester un scénario de rupture (fiche réflexe FR-07 du PCA [8, §7.2]) et la communication aux clients & Direction, achats, commercial\\
\end{xltabular}

\subsection*{Stock de sécurité, seuil de réapprovisionnement et effet sur le BFR}
Le point de commande (seuil de réapprovisionnement) vaut la consommation pendant le délai plus le stock de sécurité : $\text{PC} = \bar{D} \times L + \text{SS}$, avec $\bar{D}$ la consommation moyenne par jour et $L$ le délai en jours. Une formule courante donne $\text{SS} = z \times \sigma_D \times \sqrt{L}$, où $\sigma_D$ est l'écart-type de la consommation par jour. $z$ dépend de la probabilité visée de ne pas manquer de stock pendant le délai : environ 1,65 pour 95~\%. Ce n'est pas forcément le taux de service en volume (part de la demande servie). La formule suppose un délai fixe et une demande indépendante d'un jour à l'autre, à peu près normale. Le cas ne donne pas les consommations : le calcul se fait avec l'historique SAP.

\textit{Ordre de grandeur (illustratif).} Si la consommation de titane égale les achats du tableau (980~k€ sur 220 jours ouvrés), elle vaut environ 4,5~k€ par jour. Dix jours de stock en plus immobiliseraient environ 45~k€, soit 2~\% du stock du bilan (2\,120~k€, p.~11). Ce n'est pas un stock optimal, seulement un ordre de grandeur. Le BFR (besoin en fonds de roulement, financement du cycle d'exploitation) est déjà jugé \og un peu élevé\fg{} à 78 jours de CA (p.~37) : le stock ne doit augmenter que pour les matières critiques.

Indicateurs : taux de rupture (moins de 1~\% \tagO), délai moyen fournisseur (moins de 10 jours \tagO) [1, p.~24], part des matières critiques avec deux sources qualifiées, jours de couverture et valeur du stock de titane.
\retenir{la résilience vient surtout d'une deuxième source qualifiée pour le titane et d'un MRP fiable. Le stock de sécurité aide, mais il immobilise de la trésorerie. Il doit donc rester ciblé.}
\source{[1], p.~9--11, 24--25, 36--37 · [8], §7.2 et §9.3.}
% =====================================================================
\partie{Partie 2 -- Cybersécurité}

\question{9.2.1}{Analyse de la surface d'attaque OT/IT}{MetalTech exploite un MES Siemens communiquant en OPC-UA/MQTT avec 25 capteurs IoT et 12 machines à commande numérique, interconnecté à l'ERP. Identifiez les principaux risques de convergence IT/OT dans cette architecture et proposez trois mesures de segmentation réseau (VLAN, DMZ industrielle, firewall OT) pour limiter la propagation d'une attaque depuis le SI de gestion vers l'atelier de production.}

L'IT (technologies de l'information) est le SI de gestion : ERP, bureautique, messagerie. L'OT (technologies opérationnelles) regroupe les systèmes qui pilotent ou mesurent l'atelier : machines CN, CMM, capteurs. Chez MétalTech, le MES échange avec l'ERP par une interface de programmation (API REST) et avec les machines par OPC-UA et MQTT [1, p.~18, 20]. OPC-UA est un protocole standard d'échange de données industrielles. MQTT est un protocole de messages léger, courant pour les capteurs. Les données sont synchronisées en moins de 5 minutes (p.~18). Cette liaison est utile, mais elle ouvre un chemin d'attaque des bureaux vers l'atelier.

\subsection*{L'existant documenté}
\begin{figure}[H]
\centering
\begin{tikzpicture}[
  n/.style={box,fill=white,text width=29mm,minimum height=10mm},
  v/.style={box,fill=pale,text width=27mm,minimum height=15mm},
]
\node[n,text width=46mm] (inet) at (0,0) {Internet : 2 liens fibre\\(500 Mbit/s et 100 Mbit/s de secours)};
\node[n,text width=42mm] (az) at (-5.6,0) {Cloud Azure : réplique PRA,\\dev/test, Azure AD, sauvegardes};
\node[n,text width=42mm] (saas) at (5.6,0) {Services en ligne : Dynamics 365,\\Power BI, MES Opcenter (SaaS ?)};
\node[n,fill=pale,text width=46mm] (fg) at (0,-1.75) {\textbf{Pare-feu FortiGate 200E}\\filtrage, VPN SSL (3 utilisateurs)};
\node[n,text width=150mm,minimum height=8mm] (tor) at (0,-3.35) {Commutateurs \og Top of Rack \fg{} et segmentation en VLAN. 12 bornes Wi-Fi 6 communes à l'atelier et aux bureaux};
\node[v] (vb) at (-6.3,-5.25) {\textbf{VLAN bureau}\\45 postes fixes\\15 portables};
\node[v] (vs) at (-3.15,-5.25) {\textbf{Serveurs}\\3 serveurs HP : ERP, base HANA, PLM, fichiers, AD};
\node[v] (vp) at (0,-5.25) {\textbf{VLAN production}\\12 machines CN\\3 CMM};
\node[v] (vi) at (3.15,-5.25) {\textbf{VLAN IoT}\\25 capteurs\\(température, humidité)};
\node[v] (vg) at (6.3,-5.25) {\textbf{VLAN invités}\\Wi-Fi visiteurs};
\draw[thick,graytext] (inet)--(fg); \draw[thick,graytext] (fg)--(tor);
\draw[thick,graytext] (az)--(inet); \draw[thick,graytext] (saas)--(inet);
\foreach \x in {vb,vs,vp,vi,vg}{\draw[thick,graytext] (tor.south -| \x.north)--(\x.north);}
\node[lbl,text=alert,anchor=north] at (0,-6.35) {Règles de filtrage entre VLAN et emplacement exact des serveurs : non décrits dans le cas};
\node[lbl,text=alert,anchor=north] at (0,-6.75) {MES : SaaS selon p.~20 et 30, sur site selon p.~28 -- hébergement à confirmer};
\end{tikzpicture}
\caption{Architecture réseau existante, limitée aux éléments décrits dans le cas.}
\label{fig:itot-exist}
\figsource{D'après [1], p.~20, 26--30 et [9], §2.1 et §5.1. Aucune adresse ni règle de filtrage n'est inventée.}
\end{figure}

\needsp{12\baselineskip}
\subsection*{Principaux risques de la convergence IT/OT}
\begin{xltabular}{\linewidth}{@{}P{44mm}YP{42mm}@{}}\toprule
\thd{Risque} & \thd{Pourquoi chez MétalTech (données du cas)} & \thd{Conséquence possible}\\\midrule\endfirsthead
\toprule \thd{Risque} & \thd{Pourquoi chez MétalTech (données du cas)} & \thd{Conséquence possible}\\\midrule\endhead
\bottomrule\endfoot
Propagation d'un rançongiciel des bureaux vers l'atelier & ERP et MES synchronisés en continu (p.~18). Règles entre VLAN non décrites & Arrêt du MES et des OF, arrêt de la production\\
Accès distants de maintenance mal contrôlés & Support de niveau 2 assuré par les intégrateurs MES et SAP (p.~31). Accès par VPN SSL (p.~26) & Entrée d'un attaquant avec un compte de prestataire\\
Protocoles industriels mal configurés & OPC-UA et MQTT offrent l'authentification et le chiffrement, mais seulement s'ils sont activés & Consignes ou mesures modifiées, pièces non conformes\\
Capteurs IoT et Wi-Fi partagé & 25 capteurs. Mêmes bornes Wi-Fi pour l'atelier et les bureaux (p.~26) & Point d'entrée vers le réseau de production\\
Dépendance à Internet et au cloud & MES décrit comme un logiciel en ligne (SaaS). \og Panne internet = impact production\fg{} (p.~29) & Perte du suivi en temps réel, mode dégradé\\
Transfert des programmes CN & Programmes FAO créés côté bureaux puis envoyés aux machines (p.~5) \tagH & Programme modifié, machine endommagée, risque pour l'opérateur\\
\end{xltabular}
Dans l'atelier, la disponibilité et l'intégrité sont généralement prioritaires [2, diap.~7], sans oublier la confidentialité des plans clients. Une mauvaise consigne peut aussi mettre un opérateur en danger. Toute mesure est donc testée avec la production avant sa mise en service, et aucune règle de filtrage ne doit dégrader la sécurité des personnes.

\subsection*{Architecture cible : les trois mesures demandées}
\begin{figure}[H]
\centering
\begin{tikzpicture}[
  n/.style={box,fill=white,text width=31mm,minimum height=9mm},
  zb/.style={box,fill=pale,text width=33mm,minimum height=11mm},
  num/.style={circle,draw=teal,solid,fill=white,inner sep=0.8pt,font=\scriptsize\bfseries,text=teal},
]
\node[n,text width=150mm] (inet) at (0,0) {Internet et cloud : Azure, Dynamics 365, Power BI, MES Opcenter (si l'hébergement SaaS est confirmé)};
\node[n,fill=lightgray,text width=56mm] (fg) at (-1.6,-1.35) {Pare-feu périmétrique FortiGate (existant, mis à jour)};
% Zone IT
\node[zb] (bur) at (-5.4,-3.05) {VLAN bureaux\\postes, poste FAO};
\node[zb] (srv) at (-1.6,-3.05) {VLAN serveurs\\ERP, PLM, AD, fichiers};
\node[zb] (adm) at (2.2,-3.05) {VLAN administration\\postes d'admin SI};
\node[zb,fill=rModere!60,text width=27mm] (gst) at (6.45,-3.05) {VLAN invités\\Internet seulement,\\aucune route interne};
\begin{scope}[on background layer]
\node[zone,fit=(bur)(srv)(adm),label={[lbl,anchor=south west]north west:\textbf{Zone IT (gestion)}}] (zit) {};
\end{scope}
% Pare-feu OT et DMZ
\node[box,fill=navy,text=white,text width=58mm,minimum height=10mm] (fwot) at (-1.6,-5.3) {\textbf{Mesure 3 : pare-feu OT}\\refus par défaut, journaux};
\node[zb,fill=cB!12,draw=cB,text width=50mm,align=left] (dmz) at (4.75,-5.3) {\textbf{Mesure 2 : DMZ industrielle}\\[1pt]
\textbullet{} relais d'échange ERP--MES\\\textbullet{} dépôt des programmes CN\\\textbullet{} bastion d'administration\\\textbullet{} relais de mises à jour et d'antivirus};
% Zone OT
\node[zb] (mesl) at (-5.4,-7.95) {VLAN MES local\\connecteur MES, collecte OPC-UA, broker MQTT};
\node[zb] (cn) at (-1.6,-7.95) {VLAN cellules CN\\12 machines};
\node[zb] (cmm) at (2.2,-7.95) {VLAN métrologie\\3 CMM, PolyWorks};
\node[zb] (iot) at (6.0,-7.95) {VLAN capteurs\\25 capteurs IoT};
\begin{scope}[on background layer]
\node[zone,draw=cA,fit=(mesl)(cn)(cmm)(iot),label={[lbl,anchor=north west,text=cA]south west:\textbf{Mesure 1 : zone OT découpée en VLAN (filtrage entre VLAN par le pare-feu OT)}}] (zot) {};
\end{scope}
% Liens
\draw[thick,graytext] (fg.north)--(fg.north |- inet.south);
\draw[thick,graytext] (fg.south)--(zit.north -| fg.south);
\draw[arr,draw=graytext] (gst.north) |- ($(fg.east)+(0,0.17)$);
\draw[arrB] (zit.south -| fwot.north)--(fwot.north) node[num,pos=0.5,right=3pt] {1} node[num,pos=0.5,right=17pt] {5} node[num,pos=0.5,right=31pt] {7};
\draw[arrB] (fwot.east)--(dmz.west);
\draw[arrB] (fwot.south)--(zot.north -| fwot.south) node[num,pos=0.5,right=3pt] {2} node[num,pos=0.5,right=17pt] {6} node[num,pos=0.5,right=31pt] {8} node[num,pos=0.5,right=45pt] {10};
\coordinate (p7) at (4.4,0);
\draw[arr,draw=teal,dashed] (dmz.north -| p7) |- node[num,pos=0.8,below=2pt] {9} ($(fg.east)+(0,-0.17)$);
\draw[arrB] (mesl.east)--(cn.west);
\node[num] at ($(mesl.north east)!0.5!(cn.north west)+(-0.19,0.24)$) {3};
\node[num] at ($(mesl.north east)!0.5!(cn.north west)+(0.19,0.24)$) {4};
% Interdit
\draw[arrX] (bur.south) -- node[pos=0.5,anchor=east,xshift=-3pt,font=\scriptsize\bfseries,text=alert,align=right] {IT $\rightarrow$ machines :\\interdit} (mesl.north);
\end{tikzpicture}
\caption{Architecture cible proposée avec les trois mesures de segmentation.}
\label{fig:itot-cible}
\figsource{Proposition du groupe, pas l'existant. Les numéros renvoient au tableau des flux, qui précise le sens de chaque connexion. Les flux 3 et 4 restent dans l'atelier, entre deux VLAN, et passent aussi par le pare-feu OT. Toute connexion entre IT et OT se termine dans la DMZ industrielle. Le pare-feu OT peut être un équipement dédié ou, en solution transitoire, des zones séparées sur un pare-feu existant.}
\end{figure}

\textbf{Mesure 1 -- VLAN.} Un VLAN (réseau local virtuel) sépare logiquement des équipements branchés sur les mêmes commutateurs [4, diap.~26, 43]. MétalTech a déjà quatre VLAN : bureau, production, invités et IoT. Nous ajoutons un VLAN serveurs et un VLAN d'administration côté gestion, et nous séparons côté atelier le MES local, les machines CN, la métrologie et les capteurs. Le passage d'un VLAN à l'autre se fait seulement par un pare-feu. Chaque réseau Wi-Fi est relié à son VLAN, et le Wi-Fi invités n'atteint aucun VLAN interne.

\textbf{Mesure 2 -- DMZ industrielle.} La DMZ (zone démilitarisée) industrielle est une zone tampon entre l'IT et l'OT [4, diap.~43--45]. Aucun poste de bureau ne se connecte directement à une machine : l'ERP et le connecteur MES de l'atelier déposent ou lisent leurs données sur un relais placé dans la DMZ. La DMZ accueille aussi le dépôt des programmes CN, le bastion (serveur sécurisé qui contrôle les accès d'administration) et les relais de mises à jour. Si la DMZ est compromise, on l'isole. L'atelier ne continue en mode local que si les programmes, les fonctions de sécurité et les autres dépendances sont sur place : c'est à vérifier machine par machine.

\textbf{Mesure 3 -- Pare-feu OT.} Il filtre tous les échanges entre IT, DMZ et OT, et entre les VLAN de l'atelier. La règle de base est le refus par défaut : seuls les flux du tableau ci-dessous sont autorisés [4, diap.~46]. Les refus sont journalisés puis analysés. Pour une PME, un pare-feu à trois interfaces (IT, DMZ, OT) suffit, après validation par un intégrateur.

\subsection*{Flux autorisés (liste fermée)}
Le tableau indique qui ouvre la connexion et comment circulent les données. HTTPS sécurise les échanges web, SFTP les transferts de fichiers. Le broker MQTT reçoit et distribue les messages.

\begin{xltabular}{\linewidth}{@{}C{6mm}P{52mm}P{32mm}Y@{}}\toprule
\thd{N°} & \thd{Connexion (client $\rightarrow$ serveur)} & \thd{Protocole (port usuel)} & \thd{Données et rôle}\\\midrule\endfirsthead
\toprule \thd{N°} & \thd{Connexion (client $\rightarrow$ serveur)} & \thd{Protocole (port usuel)} & \thd{Données et rôle}\\\midrule\endhead
\bottomrule\endfoot
1 & ERP (IT) $\rightarrow$ relais d'échange (DMZ) & HTTPS, API REST (443) & OF et gammes déposés, avancement en retour. Échange ERP--MES sans lien direct\\
2 & Connecteur MES (OT) $\rightarrow$ relais d'échange (DMZ) & HTTPS (443) & Même échange, côté atelier : OF reçus, avancement et mesures envoyés\\
3 & Connecteur MES (OT) $\rightarrow$ serveur OPC-UA de la machine ou de la CMM & OPC-UA (4840 usuel, à vérifier) & Mesures, compteurs et temps remontent dans la session\\
4 & Capteur $\rightarrow$ broker MQTT (OT), puis connecteur MES $\rightarrow$ broker & MQTT avec TLS (8883 usuel) & Le capteur publie ses mesures d'ambiance, le connecteur MES s'y abonne\\
5 & Poste FAO (IT) $\rightarrow$ dépôt des programmes (DMZ) & SFTP ou HTTPS & Dépôt d'un programme validé et signé\\
6 & Poste de la cellule CN (OT) $\rightarrow$ dépôt des programmes (DMZ) & SFTP ou HTTPS & Programme récupéré après contrôle de la signature\\
7 & Poste d'administration (IT) ou prestataire par VPN $\rightarrow$ bastion (DMZ) & HTTPS ou bureau à distance, avec MFA & Entrée unique pour les sessions d'administration\\
8 & Bastion (DMZ) $\rightarrow$ équipement OT & Bureau à distance ou SSH, session enregistrée & Maintenance ouverte sur demande, pour une durée limitée\\
9 & Relais (DMZ) $\rightarrow$ MES SaaS (Internet) & HTTPS sortant (443) & Données de production, seulement si l'hébergement SaaS est confirmé\\
10 & Postes de l'atelier (OT) $\rightarrow$ relais de mises à jour (DMZ) & HTTPS & Correctifs et signatures antivirus, sans accès direct à Internet\\
\end{xltabular}
\figsource{Ports usuels par défaut, à vérifier. MFA : authentification à plusieurs facteurs. Nous ne supposons pas que chaque machine expose un serveur OPC-UA : l'inventaire le dira. En mode \og publication-abonnement\fg{} d'OPC-UA (PubSub), la machine envoie elle-même ses données : le sens de connexion change et la règle s'adapte. Restent à inventorier : heure réseau (NTP), noms (DNS), envoi des journaux, sauvegarde des serveurs de la DMZ. Tout autre flux est refusé : bureaux ou serveurs vers machines, OT vers Internet en direct, invités vers IT ou OT, bureau à distance entre zones hors bastion.}

\subsection*{Mise en œuvre sans arrêter l'atelier}
\begin{enumerate}
\item Inventorier les équipements et les flux réels pendant deux semaines (journaux du pare-feu, observation) avant d'écrire les règles.
\item Appliquer les règles en mode \og journalisation seule\fg{}, puis bloquer pendant une fenêtre de changement (vendredi soir ou week-end, p.~30), avec un retour arrière prêt.
\item Activer l'authentification et le chiffrement OPC-UA et MQTT quand les équipements le permettent, sinon filtrer plus strictement.
\item Vérifier avec la maintenance que les arrêts d'urgence et les sécurités des machines ne dépendent pas du réseau. Tester chaque trimestre qu'un flux interdit est refusé et journalisé.
\end{enumerate}
\source{[1], p.~5, 18, 20, 26--31 · [2], diap.~7, 12 · [4], diap.~26, 37, 42--46 · [9], §2.1, §5.1.}
% ---------------------------------------------------------------------
\question{9.2.2}{Gestion des vulnérabilités et obsolescence logicielle}{Le parc de MetalTech repose sur Windows 10 (fin de support en octobre 2025) et un ERP dont le support s'arrête fin 2027, sur des équipements réseau vieux de 5 ans. En vous appuyant sur une méthodologie de gestion des risques (ex. EBIOS RM ou ISO 27005), évaluez l'impact de cette dette technique et proposez un plan de remédiation priorisé sur 18 mois.}

\subsection*{Méthode : démarche ISO 27005 simplifiée}
ISO/IEC 27005 guide la gestion des risques de sécurité de l'information en cinq temps : fixer le contexte, identifier, analyser et évaluer, traiter, suivre [5, diap.~21]. Un risque combine un actif, une menace, une vulnérabilité et un impact [5, diap.~9]. Nous notons de 1 à 4 la vraisemblance (V) et la gravité (G). Le niveau vaut $V \times G$, sur la matrice 4×4 du cours [5, diap.~11].

\textbf{Échelles.} Vraisemblance sur 18 mois : 1 peu probable, 2 possible, 3 probable, 4 très probable. Gravité pour MétalTech : 1 gêne limitée sans effet client, 2 perturbation de quelques heures rattrapable, 3 arrêt partiel, retard client ou non-conformité d'audit, 4 arrêt au-delà de la durée de reprise visée de 4~h (RTO), fuite de plans clients ou perte de certification. Seuils : 1 à 3 faible, 4 à 6 modéré, 8 et 9 élevé, 12 à 16 critique. Les notes sont des hypothèses du groupe, à valider avec le RSSI (responsable de la sécurité des SI) et la direction.

\subsection*{Appréciation des risques}
Le sujet parle d'un parc sous Windows 10, mais le cas cite 45 postes sous Windows 11 Pro (p.~28) et une migration prévue au premier trimestre 2026 (p.~33). La première action est donc un inventaire des postes et de leurs éditions. L'âge d'un équipement n'est pas une vulnérabilité en soi : c'est l'absence de correctifs ou de support qui en crée une.

\begin{xltabular}{\linewidth}{@{}C{7mm}P{28mm}YC{15mm}YC{16mm}@{}}\toprule
\thd{Id} & \thd{Actif et menace} & \thd{Vulnérabilité (donnée du cas)} & \thd{Initial V×G} & \thd{Mesures principales} & \thd{Résiduel estimé}\\\midrule\endfirsthead
\toprule \thd{Id} & \thd{Actif et menace} & \thd{Vulnérabilité (donnée du cas)} & \thd{Initial V×G} & \thd{Mesures principales} & \thd{Résiduel estimé}\\\midrule\endhead
\bottomrule\endfoot
R1 & Postes encore sous Windows 10. Rançongiciel, hameçonnage & Plus de correctifs standard pour la version 22H2 depuis le 14/10/2025 [18]. Nombre réel inconnu & 4×4 = \textbf{16} & Inventaire, isolement, mises à jour de sécurité étendues (ESU) en transition, migration vers Windows 11 & 1×4 = 4\\
R2 & ERP SAP Business One 10.0. Arrêt de la gestion et de la production & Fin de support annoncée fin 2027 dans le cas (p.~33). ERP vital, priorité P1 [9, §2.1] & 3×4 = \textbf{12} & Choisir une trajectoire. Migration ou support transitoire validé avant fin 2027 & 1×4 = 4\\
R3 & Pare-feu FortiGate 200E et commutateurs. Intrusion, panne & 5 ans d'âge (p.~33). Pare-feu unique et VPN exposé à Internet & 3×4 = \textbf{12} & Firmware à jour, configuration sauvegardée, MFA sur le VPN, remplacement planifié & 1×4 = 4\\
R4 & 12 postes de plus de 5 ans sur 60. Panne, incompatibilité & Matériel ancien (p.~33), peut-être incompatible avec Windows 11 & 3×2 = 6 & Remplacement dans le plan matériel de 25~k€ par an & 1×2 = 2\\
R5 & Interfaces SAP--MES--PLM pendant la migration. Erreurs de données & Flux en temps réel (p.~18), personnalisations de l'ERP (p.~19) & 3×3 = 9 & Pilote, tests des interfaces, retour arrière préparé & 2×2 = 4\\
R6 & Équipe SI réduite. Correctifs en retard & 2 personnes en poste et 1 poste vacant (p.~30). Pas de SOC, seulement un projet (p.~33) & 3×3 = 9 & Calendrier de correctifs, appui d'un prestataire, SOC externalisé & 2×3 = 6\\
\end{xltabular}
\figsource{Le risque résiduel est celui qui reste après les mesures de sécurité. Ici, il est estimé. ESU : mises à jour de sécurité étendues, payantes. SOC : centre de supervision de la sécurité.}
\begin{figure}[htp]
\centering
\begin{tikzpicture}[c/.style={minimum size=10mm,inner sep=0pt,draw=white,line width=1pt,font=\scriptsize,align=center}]
\newcommand{\riskcolor}[2]{\pgfmathtruncatemacro{\s}{#1*#2}%
  \ifnum\s<4 \def\cc{rFaible}\else\ifnum\s<7 \def\cc{rModere}\else\ifnum\s<10 \def\cc{rEleve}\else\def\cc{rCritique}\fi\fi\fi}
\foreach \ox/\titre in {0/{Avant traitement},6.9/{Après traitement (estimé)}}{
  \begin{scope}[xshift=\ox cm]
  \foreach \v in {1,...,4}{\foreach \g in {1,...,4}{
     \riskcolor{\v}{\g}
     \node[c,fill=\cc] at (\g,\v) {};}}
  \foreach \g in {1,...,4}{\node[font=\scriptsize] at (\g,0.3) {G\g};}
  \foreach \v in {1,...,4}{\node[font=\scriptsize] at (0.3,\v) {V\v};}
  \node[font=\footnotesize\bfseries,text=navy] at (2.5,4.85) {\titre};
  \node[font=\scriptsize] at (2.5,-0.2) {Gravité};
  \node[font=\scriptsize,rotate=90] at (-0.25,2.5) {Vraisemblance};
  \end{scope}}
% positions initiales
\node[font=\scriptsize\bfseries] at (4,4) {R1};
\node[font=\scriptsize\bfseries,align=center] at (4,3) {R2\\R3};
\node[font=\scriptsize\bfseries] at (2,3) {R4};
\node[font=\scriptsize\bfseries,align=center] at (3,3) {R5\\R6};
% positions résiduelles
\node[font=\scriptsize\bfseries,align=center] at (6.9+4,1) {R1 R2\\R3};
\node[font=\scriptsize\bfseries] at (6.9+2,1) {R4};
\node[font=\scriptsize\bfseries] at (6.9+2,2) {R5};
\node[font=\scriptsize\bfseries] at (6.9+3,2) {R6};
% légende
\foreach \c/\t [count=\k] in {rFaible/{faible (1--3)},rModere/{modéré (4--6)},rEleve/{élevé (8--9)},rCritique/{critique (12--16)}}{
  \node[fill=\c,minimum size=3.5mm,inner sep=0pt] at (12.6,{4.3-0.6*\k}) {};
  \node[font=\scriptsize,anchor=west] at (12.85,{4.3-0.6*\k}) {\t};}
\end{tikzpicture}
\caption{Matrice des risques 4×4 avant et après traitement.}
\label{fig:risques}
\figsource{Notes : hypothèses du groupe. Le risque résiduel est une estimation à confirmer après la mise en œuvre des mesures. La direction le valide formellement [5, diap.~12, 21].}
\end{figure}

\subsection*{Impact de la dette technique}
Selon les seuils de la matrice, R1 (16), R2 et R3 (12) sont des risques critiques sur des actifs vitaux. L'ERP, le réseau local et le pare-feu sont classés P1, avec un RTO (durée maximale visée avant la reprise) de 4~h [9, §4]. Un rançongiciel entré par un poste non maintenu ou par le VPN peut arrêter la gestion et l'atelier. En ordre de grandeur, MétalTech réalise environ 84~k€ de CA par jour ouvré (18,4~M€ / 220 jours, hypothèse). Ce n'est pas une perte certaine : le coût réel d'un arrêt dépend de la marge, des pénalités et du rattrapage possible. La dette technique freine aussi l'ISO 27001, car un auditeur attend une gestion maîtrisée des vulnérabilités.

\subsection*{Plan de remédiation priorisé sur 18 mois}
Nous fixons M0 à décembre 2026. M12 correspond donc à décembre 2027, échéance du support de l'ERP selon le sujet. M18 correspond à juin 2028.
\begin{figure}[htp]
\centering
\begin{tikzpicture}[x=0.5cm,y=0.64cm]
\def\L{-0.3}
% grille des mois
\foreach \m in {0,...,18}{\draw[lightgray] (\m,0.6)--(\m,-12.4);}
\foreach \m in {0,3,...,18}{\draw[midgray] (\m,0.6)--(\m,-12.4); \node[font=\scriptsize\bfseries] at (\m,1.05) {M\m};}
\foreach \m/\t in {0/déc. 26,3/mars 27,6/juin 27,9/sept. 27,12/déc. 27,15/mars 28,18/juin 28}{\node[font=\tiny,text=graytext] at (\m,1.6) {\t};}
% tâches : ligne, début, fin, couleur, libellé, responsable
\foreach \r/\a/\b/\c/\t/\p in {
  1/0/1/cA/{Inventaire des postes, versions et firmwares, scan}/{SI, RSSI},
  2/0/3/cA/{Isoler ou protéger les postes non supportés (ESU)}/{SI},
  3/1/5/cA/{Migrer vers Windows 11, remplacer les 12 postes}/{SI, métiers},
  4/0/2/cB/{Mettre à jour FortiGate et commutateurs}/{SI, intégrateur},
  5/1/3/cB/{MFA sur le VPN et durcissement (voir 9.2.6)}/{SI},
  6/3/8/cB/{Remplacer pare-feu et commutateurs prioritaires}/{SI, intégrateur},
  7/1/3/cC/{Étudier les options ERP et le budget}/{DSI, DAF},
  8/4/9/cC/{Pilote, reprise des données, tests SAP--MES--PLM}/{SI, intégrateur SAP},
  9/9/12/cC/{Bascule phase 1 ou support transitoire validé}/{DSI, direction},
  10/12/18/cC/{Stabiliser, traiter les lots restants, retirer l'ancien}/{DSI, métiers},
  11/2/18/cD/{Correctifs mensuels et scans de vulnérabilités}/{SI, SOC (projet)}}{
  \fill[\c,rounded corners=1.5pt] (\a,{-\r+0.28}) rectangle (\b,{-\r-0.28});
  \node[font=\scriptsize,anchor=east,align=right,inner sep=1pt] at (\L,-\r) {\t\\[-3pt]{\tiny\textcolor{graytext}{\p}}};
}
\node[font=\scriptsize,anchor=east,align=right,inner sep=1pt] at (\L,-12) {Revues du risque résiduel par la direction\\[-3pt]{\tiny\textcolor{graytext}{RSSI, direction}}};
\foreach \m in {6,12,18}{\node[diamond,fill=cD,inner sep=2.2pt] at (\m,-12) {};}
% jalons
\node[diamond,fill=navy,inner sep=2.4pt] at (3,-7) {};
\node[font=\tiny,anchor=west] at (3.2,-6.55) {décision};
\node[diamond,fill=alert,inner sep=2.6pt] at (12,-9) {};
\node[font=\tiny,text=alert,anchor=west,align=left] at (12.25,-8.35) {support ERP assuré\\(fin 2027 selon le sujet)};
% dépendances
\draw[arr,thin] (1,-1)--(1,-3);
\draw[arr,thin] (3,-7)--(4,-8);
\draw[arr,thin] (9,-8)--(9,-9);
\end{tikzpicture}
\caption{Calendrier de remédiation sur 18 mois (M0 = décembre 2026).}
\label{fig:gantt}
\figsource{Proposition du groupe. Couleurs : bleu pour les postes, vert-bleu pour le réseau, orange pour l'ERP, violet pour le pilotage. Le libellé de chaque tâche permet aussi la lecture en niveaux de gris. Flèches : dépendances principales. Losanges : jalons.}
\end{figure}

Les coûts du plan (matériel, ERP, SOC, ESU, pare-feu) sont regroupés une seule fois dans le budget de la conclusion, avec leur statut et leur source.

\textbf{Faisabilité avec trois personnes.} Le plan est chargé pour deux personnes en poste et un poste vacant. Trois règles : ne pas mener la migration Windows et le pilote ERP sur les mêmes semaines ; confier l'ERP et le pare-feu à des intégrateurs, l'équipe SI pilotant et validant ; pourvoir le poste vacant avant le pilote ERP (M4).

\textbf{Vérifications externes (24 et 25/11/2026).} Le support standard de Windows 10 version 22H2 a pris fin le 14 octobre 2025. Certaines éditions LTSC (support long) suivent un autre calendrier, par exemple Entreprise LTSC 2021 jusqu'au 12 janvier 2027 [18] : l'inventaire relève donc l'édition de chaque poste. Des ESU payantes permettent une transition [16]. Pour SAP Business One 10.0, SAP indique une maintenance principale jusqu'au 31 décembre 2028, date \og susceptible de changer\fg{} [17]. Nous gardons l'échéance de fin 2027 du sujet, à confirmer avec l'intégrateur. Le cycle de vie du FortiGate 200E et sa version de logiciel interne (firmware) n'ont pas pu être vérifiés sur une source officielle publique : ils sont à contrôler sur le portail Fortinet. Nous ne citons aucune faille sans preuve.

Indicateurs mensuels du comité SI : part des postes sous un système supporté (cible 100~\%), délai de correction des failles critiques, exceptions approuvées avec date de fin, avancement des tests de migration de l'ERP.
\source{[1], p.~18--19, 26--33 · [5], diap.~9--12, 21 · [9], §2.1, §4 · [16] · [17] · [18].}
% ---------------------------------------------------------------------
\question{9.2.3}{Protection de la propriété intellectuelle sensible (PLM/CAO)}{Le PLM Windchill héberge 120 Go de données CAO incluant des spécifications liées à l'aéronautique (A320). Quelles mesures techniques et organisationnelles (chiffrement, DLP, gestion des accès, traçabilité) recommanderiez-vous pour protéger ces données contre l'espionnage industriel, en tenant compte des exigences EN 9100/NADCAP ?}

Le PLM PTC Windchill contient 120~Go de données de CAO, 2\,800 articles et 450 dossiers de définition [1, p.~21]. Les clients échangent des plans par un portail Windchill (STEP, IGES, PDF 3D). Le cas cite des attaques par hameçonnage ciblé sur les données techniques (p.~33). Les plans A320 doivent rester confidentiels et intègres : une cote modifiée peut produire une pièce non conforme.

\subsection*{Classer les données avant de les protéger}
\begin{xltabular}{\linewidth}{@{}P{30mm}YY@{}}\toprule
\thd{Niveau proposé} & \thd{Exemples chez MétalTech} & \thd{Règle principale}\\\midrule\endfirsthead
\toprule \thd{Niveau proposé} & \thd{Exemples chez MétalTech} & \thd{Règle principale}\\\midrule\endhead
\bottomrule\endfoot
Public & Plaquette, références générales & Diffusion libre\\
Interne & Procédures, gammes standard & Accès pour les salariés, pas de diffusion externe\\
Confidentiel & Devis, coûts, nomenclatures & Accès par rôle, partage externe validé\\
Confidentiel client & Plans et spécifications A320, données fournies par les clients & Accès nominatif par projet, MFA, traçabilité complète, partage seulement par le portail\\
\end{xltabular}
\figsource{Proposition du groupe. Chaque dossier a un propriétaire, le responsable du bureau d'études ou du projet, qui valide les accès.}

\subsection*{Quelques notions simples}
\begin{itemize}
\item \textbf{AES} est un chiffrement symétrique (la même clé chiffre et déchiffre), utilisé pour les données stockées [3, diap.~15]. \textbf{TLS} chiffre les échanges sur le réseau, par exemple entre le portail et un client [3, diap.~25].
\item Une \textbf{empreinte} (hachage) détecte la modification d'un fichier [3, diap.~21]. Si un attaquant peut remplacer le fichier et son empreinte, le contrôle ne sert à rien : les programmes CN et les plans sensibles demandent une \textbf{signature numérique}, ou une empreinte transmise par un canal de confiance.
\item En résumé : le chiffrement protège la \emph{confidentialité}, le hachage contrôle l'\emph{intégrité}, AES-GCM (chiffrement authentifié) protège les deux [19], et une signature vérifiée avec une clé de confiance prouve l'origine du fichier, son \emph{authenticité}.
\item Le \textbf{DLP} (prévention des fuites de données) détecte, signale ou bloque les transferts non autorisés de données sensibles. Le \textbf{RBAC} (contrôle d'accès par rôle) donne les droits selon le rôle et le projet [6, diap.~30].
\end{itemize}

\subsection*{Mesures techniques et organisationnelles}
\begin{xltabular}{\linewidth}{@{}P{27mm}YP{38mm}@{}}\toprule
\thd{Mesure} & \thd{Application à MétalTech} & \thd{Preuve ou contrôle}\\\midrule\endfirsthead
\toprule \thd{Mesure} & \thd{Application à MétalTech} & \thd{Preuve ou contrôle}\\\midrule\endhead
\bottomrule\endfoot
Chiffrement au repos & BitLocker, déjà en place (p.~26), protège contre le vol d'un disque ou d'un ordinateur, mais pas contre un utilisateur autorisé ni un rançongiciel actif dans une session ouverte. Ajouter, si Windchill le permet, le chiffrement du coffre de fichiers et de la base (AES-256) [3, diap.~15, 39] & Configuration documentée, test de protection des données et des clés\\
Chiffrement en transit & TLS 1.2 au minimum, 1.3 de préférence, pour le portail et les échanges [3, diap.~25] & Test de configuration TLS\\
Gestion des clés & Clés séparées des données, accès réservé à deux administrateurs, journal des accès [3, diap.~39, 42] & Inventaire des clés, revue annuelle\\
DLP & Alerter puis bloquer la copie de fichiers CAO vers une clé USB, une messagerie personnelle ou un cloud non autorisé, d'abord avec les licences existantes & Journaux d'alertes, exceptions validées\\
Gestion des accès & Comptes nominatifs, droits par rôle et par projet (RBAC), MFA pour l'accès distant et le portail, revue trimestrielle [6, diap.~6, 30, 32] & Matrice des accès, compte rendu de revue\\
Partenaires & Accès au portail limité au projet et dans le temps, accord de confidentialité, filigrane sur les PDF exportés & Liste à jour des comptes externes\\
Traçabilité & Journaux Windchill (consultations, téléchargements, modifications), alerte sur les téléchargements massifs & Revue des alertes\\
Conservation des journaux & Durée fixée avec la qualité et le DPO (délégué à la protection des données). Copie hors du serveur, protégée contre la modification, lisible par le RSSI et l'équipe SI & Procédure écrite, contrôle des accès\\
Intégrité et versions & Validation et versions dans Windchill (p.~21). Signature numérique, ou empreinte transmise par un canal de confiance, des programmes et plans envoyés aux machines et aux clients [3, diap.~21] & Contrôle de signature, historique des versions\\
Sauvegardes & Sauvegarde séparée du PLM [8, §3.4], copie hors ligne ou immuable, test de restauration de la bonne version & Rapport de test\\
Facteur humain & Sensibilisation à l'hameçonnage ciblé, d'abord au bureau d'études et chez les commerciaux [2, diap.~25]. Révocation des accès le jour du départ [6, diap.~35] & Taux de clic, délai de révocation\\
\end{xltabular}
Pour les données qui doivent rester secrètes longtemps, le cours signale le risque \og stocker maintenant, déchiffrer plus tard\fg{} [3, diap.~33] : il faut pouvoir changer d'algorithme plus tard (agilité cryptographique) [3, diap.~36].

\subsection*{Lien avec EN 9100 et NADCAP}
EN 9100 est un référentiel qualité pour l'aéronautique. NADCAP accrédite les procédés spéciaux, ici les traitements thermiques [1, p.~16, 25]. Ils demandent surtout des documents maîtrisés, des versions à jour, la traçabilité des produits et des enregistrements fiables. Nos mesures les soutiennent : versions et validations dans Windchill, signature des fichiers, sauvegardes testées, traçabilité des accès. Nous n'affirmons pas qu'EN 9100 ou NADCAP imposent un algorithme ou un outil DLP précis : les exigences de confidentialité viennent plutôt des contrats clients, à vérifier.

\textbf{Priorités.} M0--M1 : classer les dossiers A320, nommer les propriétaires, revoir les comptes du portail (bureau d'études, RSSI). M1--M3 : MFA pour le portail et les accès distants, TLS vérifié, sauvegarde séparée testée (SI). M3--M6 : DLP en alerte puis en blocage, journaux centralisés, première revue trimestrielle des accès (SI, qualité).
\source{[1], p.~16, 21, 25--27, 33 · [2], diap.~25 · [3], diap.~15, 21, 25, 33, 36, 39, 42 · [6], diap.~6, 30, 32, 35 · [8], §3.4 · [19].}
% ---------------------------------------------------------------------
\question{9.2.4}{Réponse à incident : scénario ransomware}{En supposant un scénario de ransomware ciblant l'ERP et les serveurs de fichiers de MetalTech (PME industrielle, équipe SI de 3 personnes, absence de SOC), rédigez un plan de réponse à incident simplifié (détection, confinement, éradication, restauration via Veeam/Azure, communication) adapté aux ressources limitées de l'entreprise.}

\subsection*{Hypothèses et moyens disponibles}
Selon le sujet, un rançongiciel (logiciel qui chiffre les données et demande une rançon) touche l'ERP et les serveurs de fichiers, avec une équipe SI de 3 personnes et sans SOC. Dans le cas, l'équipe compte un responsable SI, un administrateur SI et un poste vacant (p.~30), plus un RSSI externe à mi-temps (p.~13). Le SOC externalisé (environ 15~k€ par an) n'est qu'un projet (p.~33).

Moyens en place [1, p.~26--27, 31] : Sophos Intercept X sur les postes et les serveurs (fonctions EDR, détection et réponse sur les postes, à vérifier dans la licence), pare-feu FortiGate, sauvegardes Veeam (instantanés toutes les 4~h, gardés 30 jours), réplique des serveurs sur Azure, archives Azure Blob, outil de tickets Jira, cyberassurance AXA (2~M€) et DPO externe. La réplication vers Azure est asynchrone, environ toutes les 15 minutes [9, annexe B]. Une réplique n'est pas une sauvegarde immuable : si elle reste connectée, le rançongiciel peut l'atteindre.

\subsection*{Déclenchement : qui décide et quand}
Nous reprenons les niveaux du cours [7, diap.~12, 16] :
\begin{itemize}
\item \textbf{Niveau 1, incident} : alerte isolée sur un poste. L'équipe SI traite et documente dans Jira.
\item \textbf{Niveau 2, incident majeur} : plusieurs alertes, fichiers renommés ou chiffrés sur un partage, compte à privilèges suspect. Le responsable SI et le RSSI pilotent, la direction est informée.
\item \textbf{Niveau 3, crise} : ERP ou serveurs de fichiers chiffrés, note de rançon, production touchée ou vol de données probable. Le PDG active la cellule de crise, le PCA (plan de continuité d'activité) et le PAR (plan d'assurance des reprises, partie technique du PCA, appelée PRA dans le nom du fichier). Dans notre scénario, ce niveau est atteint dès que l'ERP est chiffré.
\end{itemize}
\begin{figure}[htp]
\centering
\begin{tikzpicture}[x=0.8cm,y=0.62cm]
% axe non linéaire : positions des repères
\foreach \x/\t in {0/H0,1.4/H+1,2.8/H+4,4.6/H+24,6.4/H+72,8.4/J+5,10.2/J+10,12.6/J+30}{
  \draw[lightgray] (\x,0.6)--(\x,-9.2); \node[font=\scriptsize\bfseries] at (\x,1.0) {\t};}
\foreach \r/\a/\b/\c/\t in {
  1/0/1.4/cA/{Détection, qualification, journal de crise},
  2/0.3/2.8/cA/{Confinement : isoler, couper les flux},
  3/0.3/4.6/cD/{Conservation des preuves},
  4/1.4/8.4/cB/{Mode dégradé de l'atelier (PCA)},
  5/2.8/8.4/cC/{Investigation et éradication},
  6/4.6/8.4/cC/{Restauration propre en zone isolée},
  7/8.4/12.6/cB/{Validation métier, retour sous surveillance}}{
  \fill[\c,rounded corners=1.5pt] (\a,{-\r+0.32}) rectangle (\b,{-\r-0.32});
  \node[font=\scriptsize,anchor=east] at (-0.2,-\r) {\t};}
% notifications
\node[font=\scriptsize,anchor=east] at (-0.2,-8.2) {Notifications et RETEX};
\foreach \x/\t/\dy/\an in {4.6/{assureur, alerte\\précoce NIS2}/0.4/north, 6.4/{notification NIS2,\\CNIL si risque}/1.5/north, 10.2/{RETEX}/0.4/north, 13.05/{rapport final NIS2 :\\1 mois après la notification}/1.5/north east}{
  \node[diamond,fill=alert,inner sep=2.2pt] at (\x,-8.2) {};
  \draw[alert,thin] (\x,-8.2)--(\x,{-8.2-\dy+0.15});
  \node[font=\tiny,align=center,anchor=\an,fill=white,inner sep=1pt] at (\x+0.15,{-8.2-\dy+0.15}) {\t};}
\end{tikzpicture}
\caption{Chronologie indicative de la réponse au rançongiciel.}
\label{fig:crise}
\figsource{Proposition du groupe. Axe non proportionnel. Délais indicatifs et non testés, identiques au tableau ci-dessous : restauration de J+1 à J+5, puis validation métier et retour progressif jusqu'à J+30 au plus tard. Le RTO de 4~h est alors dépassé : recours au mode dégradé du PCA. H0 est le moment où l'entreprise a connaissance de l'incident. Les délais NIS2 partent de ce moment, sauf le rapport final, dû un mois après la notification (vers J+33 si la notification part à H+72). Voir 9.2.5.}
\end{figure}

\subsection*{Plan de réponse opérationnel}
\begin{xltabular}{\linewidth}{@{}P{28mm}YP{22mm}@{}}\toprule
\thd{Phase} & \thd{Actions concrètes} & \thd{Délai indicatif}\\\midrule\endfirsthead
\toprule \thd{Phase} & \thd{Actions concrètes} & \thd{Délai indicatif}\\\midrule\endhead
\bottomrule\endfoot
Détection et qualification & Repérer les alertes Sophos, les fichiers renommés ou l'ERP indisponible. Noter l'heure, les systèmes touchés et les indices dans un journal de crise tenu hors du SI. Prévenir le responsable SI, le RSSI et le PDG & H0 à H+1\\
Sécurité des personnes & En cas de danger immédiat, la mise en sécurité des personnes passe avant tout. Sinon, l'arrêt contrôlé des machines est décidé avec le chef d'atelier & Dès H0\\
Confinement sûr & Isoler les machines touchées depuis la console EDR si disponible, ou en débranchant le réseau, sans les éteindre : la mémoire contient des preuves [7, diap.~19]. Couper la liaison Azure [9, §7.3]. Séparer IT et OT avec les règles FortiGate et les VLAN existants (le pare-feu OT de 9.2.1 n'est pas encore en place). Désactiver les comptes suspects, isoler le serveur Veeam & Dès H0, sans attendre le diagnostic\\
Preuves & Garder les journaux FortiGate, Sophos, Active Directory et Azure, photographier la note de rançon, ne rien effacer tant que la sécurité le permet. Prévenir l'assureur, qui peut fournir un prestataire (à vérifier dans le contrat). Préparer la plainte & H0 à H+24\\
Éradication & Trouver la porte d'entrée (hameçonnage, VPN, poste non à jour) et la date réelle de l'intrusion. Changer les mots de passe ne suffit pas : assainir ou reconstruire l'annuaire, révoquer sessions et jetons, renouveler secrets et comptes à privilèges, corriger la faille, reconstruire depuis des sources sûres [7, diap.~20] & H+4 à J+5\\
Restauration propre & Restaurer dans une zone isolée (VLAN dédié, PAR §8.1) depuis une sauvegarde antérieure à l'intrusion, pas seulement au chiffrement. L'analyser avant la remise en service. Ne pas reconnecter Azure automatiquement : la réplique peut contenir des données chiffrées ou des traces de l'intrusion & J+1 à J+5\\
Validation et retour & Chaque service est validé par le métier (production, qualité, ADV) avant sa réouverture. Ressaisir les données du mode dégradé, contrôler les doublons. Surveillance renforcée au moins 48~h [9, §5.3] & J+5 à J+30\\
Communication et RETEX & Points de situation sur un canal hors SI. Notifications légales (9.2.5). Retour d'expérience (RETEX) sous 10 jours [8, §7.1], bilan à J+30 [7, diap.~25], mise à jour du PCA et du PAR & Toute la crise\\
\end{xltabular}
Nous ne recommandons pas de payer la rançon : le paiement ne garantit ni le déchiffrement ni la suppression des données volées. Une restauration propre peut durer plusieurs jours et dépasser le RTO de 4~h. L'activité repose alors sur le mode dégradé du PCA, et la direction choisit les commandes prioritaires, les clients à prévenir et la part éventuelle de sous-traitance.

\subsection*{Responsabilités (RACI)}
La matrice RACI indique qui fait quoi [7, diap.~13] : R réalise, A approuve (seul responsable final), C est consulté, I est informé. Chaque ligne a un seul A et au moins un R.

\par\noindent{\footnotesize\begin{tabularx}{\linewidth}{@{}YC{8.5mm}C{10mm}C{10.8mm}C{10mm}C{10.5mm}C{12mm}C{10mm}C{10mm}C{11.5mm}@{}}\toprule
\thd{Action} & \thd{PDG} & \thd{Resp. SI} & \thd{Admin SI} & \thd{RSSI ext.} & \thd{Dir. prod.} & \thd{Qualité} & \thd{Dir. com.} & \thd{DPO ext.} & \thd{Presta\-taire}\\\midrule
Activer la cellule de crise (niveau 3) & A & R & I & C & I & I & I & I & --\\
Qualifier l'incident, tenir le journal & I & A & R & C & I & -- & -- & I & C\\
Confinement du réseau et des comptes & I & A & R & C & C & -- & -- & -- & R\\
Mode dégradé de l'atelier & I & C & -- & -- & A & R & I & -- & --\\
Conservation des preuves & I & R & R & A & -- & -- & -- & I & C\\
Éradication et reconstruction & I & A & R & C & -- & -- & -- & -- & R\\
Restauration et ordre de reprise & I & A & R & C & C & C & -- & -- & R\\
Autoriser le retour en production & A & C & -- & C & R & R & I & -- & --\\
Notifications NIS2, CNIL, assureur & A & C & -- & R & -- & -- & -- & R & --\\
Communication aux clients et aux salariés & A & C & -- & C & C & C & R & C & --\\
RETEX et plan d'actions & A & R & C & R & C & C & C & C & C\\\bottomrule
\end{tabularx}}
\figsource{Le directeur de crise est le PDG, avec la direction de production comme suppléance [8, §6.1]. La directrice commerciale porte la communication, comme le prévoit le PCA [8, §6.1, §7.3]. Prestataire : intégrateur ou spécialiste de la réponse à incident.}

\subsection*{Continuité, reprise et écarts à traiter}
Le PCA maintient les activités essentielles en mode dégradé, le PAR organise la reprise technique du SI. Le RTO est la durée maximale visée avant la reprise, le RPO la durée maximale de données que l'on accepte de perdre [7, diap.~22].

\begin{xltabular}{\linewidth}{@{}P{40mm}C{11mm}C{11mm}YP{33mm}@{}}\toprule
\thd{Activité ou service} & \thd{RTO} & \thd{RPO} & \thd{Mode dégradé prévu} & \thd{Écart ou point à vérifier}\\\midrule\endfirsthead
\toprule \thd{Activité ou service} & \thd{RTO} & \thd{RPO} & \thd{Mode dégradé prévu} & \thd{Écart ou point à vérifier}\\\midrule\endhead
\bottomrule\endfoot
Production (PCA, ACT-01) & 4 h & 1 h & OF papier, relevés manuels & Dépend de l'ERP, du MES et du PLM\\
ERP et base HANA (PAR, TIC-01 et 02) & 4 h & 1 h & Réplique Azure, instantanés Veeam & Instantanés toutes les 4~h. Fréquence des journaux de la base HANA \og à confirmer\fg{} [9, §6.1]\\
MES Opcenter (PAR, TIC-03) & 4 h & 4 h & Suivi manuel, OF papier & RPO de 4~h plus large que la cible de 1~h de la production\\
PLM Windchill (PAR, TIC-04) & 8 h & 8 h & Cache local de CAO, pas de nouveaux OF & Plans nécessaires pour lancer de nouvelles séries\\
Qualité (PCA, ACT-02) & 8 h & 4 h & Check-lists papier, PolyWorks local & Traçabilité à reconstituer après la reprise\\
\end{xltabular}
\figsource{Valeurs vérifiées dans les originaux [8, §4.1] et [9, §4]. Ce sont des objectifs, pas des garanties. Les documents présentent eux-mêmes plusieurs valeurs comme des estimations à valider en atelier BIA (analyse d'impact sur l'activité).}

Trois écarts sont essentiels. Le test de mars 2026 (RTO 4~h, RPO 1~h) [1, p.~27] portait sur une bascule de secours : il ne garantit pas ces délais après un rançongiciel. Des instantanés toutes les 4~h ne suffisent pas, seuls, pour un RPO de 1~h : il faut vérifier les journaux de transactions de la base et les sauvegardes de chaque application. Enfin, une copie Azure connectée peut être atteinte. Nous proposons la règle 3-2-1 : trois copies, sur deux supports, dont une hors site [7, diap.~10], et l'une immuable ou hors ligne. Le PAR prévoit déjà un test de restauration par trimestre [9, §6.3].

\subsection*{Ordre de reprise proposé}
\begin{enumerate}
\item Préparer une zone propre et isolée, vérifier le serveur de sauvegarde, choisir un point de restauration sain.
\item Rétablir réseau et pare-feu, puis l'annuaire (Active Directory et Entra ID, ancien Azure AD) : assaini ou reconstruit, sessions révoquées, comptes à privilèges renouvelés.
\item Restaurer la base HANA, l'ERP et la plateforme d'intégration (services P1 : 4~h et 1~h), validés par l'ADV et la production.
\item Restaurer les fichiers critiques, d'abord ceux de la production et de la qualité.
\item Reconnecter le MES, puis l'atelier, après contrôle des OF et des lots par la qualité.
\item Restaurer et valider le PLM (ou un cache de plans autorisé) et la traçabilité qualité. Vérifier l'intégrité des plans (versions validées, signatures ou empreintes) avant toute nouvelle série qui en dépend.
\item Rétablir ensuite les services P2 et P3.
\end{enumerate}
Pour les séries déjà lancées, le mode dégradé reste possible si le programme et le plan validés sont sur place, si les contrôles qualité peuvent être faits et enregistrés, et si les fonctions de sécurité des machines restent opérationnelles.

\subsection*{Communication et préparation}
Le PCA prévoit un canal Teams pour la cellule de crise [8, §7.1], mais Teams dépend de comptes qui peuvent être compromis. Nous proposons un canal de secours préparé : liste papier des contacts (annexe B du PCA), téléphones mobiles, messagerie externe [7, diap.~10]. Une seule personne parle à l'extérieur, avec des messages factuels et datés [7, diap.~27]. Le PCA prévoit de prévenir les clients clés touchés dans les 4~h (12 clients aéronautiques et 20 principaux clients automobiles) et l'assureur dans les 24~h [8, §7.3] : ce sont des délais internes, pas des délais légaux.

À faire dès maintenant, avec peu de moyens : imprimer la fiche réflexe rançongiciel (FR-03 du PCA) et la liste des contacts, créer et tester des comptes d'urgence et une copie hors ligne, vérifier l'assistance prévue par le contrat AXA, faire un exercice de crise par an et un test de restauration par trimestre [8, §8.1]. Le SOC externalisé surveillerait les alertes la nuit et le week-end (budget en conclusion).
\source{[1], p.~13, 26--27, 30--33 · [7], diap.~10, 12--13, 16, 18--20, 22, 25, 27 · [8], §4.1, §6.1, §7.1--7.3, §8.1 · [9], §4, §5.3, §6, §7.3, §8.1, annexe B.}
% ---------------------------------------------------------------------
\question{9.2.5}{Conformité réglementaire (RGPD / NIS2)}{MetalTech est qualifiée d'\og entité importante\fg{} au sens de la directive NIS2 et traite des données personnelles soumises au RGPD. Quelles obligations concrètes (notification d'incident, mesures de sécurité minimales, gouvernance) découlent de ce double cadre réglementaire, et quelles actions prioritaires proposeriez-vous pour atteindre la certification ISO 27001 visée au T4 2026 ?}

\subsection*{Cadre retenu et état de la transposition}
Nous suivons l'hypothèse du sujet : MétalTech est une \og entité importante\fg{} au sens de NIS2. Le cas dit ce statut \og en cours d'évaluation\fg{} (p.~27) ou \og probable\fg{} (p.~34). Il dépend en réalité du secteur (code NACE) et de la taille : l'annexe II de la directive vise par exemple la construction automobile et les autres matériels de transport (divisions 29 et 30 de la NACE) [10]. Une pièce métallique peut relever d'une autre division : à vérifier.

Quatre sources sont à distinguer : la \textbf{directive NIS2} (UE) 2022/2555, qui fixe des exigences minimales aux États [10] ; la \textbf{loi française de transposition}, qui fixera les obligations applicables en France ; le \textbf{RGPD} (UE) 2016/679, d'application directe [11] ; les \textbf{recommandations de l'ANSSI} (Agence nationale de la sécurité des systèmes d'information), utiles mais non obligatoires [13].

\textbf{Vérification au 25 novembre 2026.} L'Assemblée nationale a adopté en première lecture, le 17 novembre 2026, le projet de loi relatif à la résilience des infrastructures critiques et au renforcement de la cybersécurité, en modifiant le texte voté par le Sénat le 12 mars 2025. La procédure accélérée est engagée et le dossier ne mentionne aucune étape après le 17 novembre (mise à jour du 18/11/2026) [14] : le texte n'est pas promulgué. MesServicesCyber (ANSSI) indique une transposition en cours, avec un pré-enregistrement des entités et le référentiel ReCyF [13] : une démarche utile, pas encore obligatoire. Le cours présente la \og loi n°~2025-379\fg{} comme loi de transposition [2, diap.~28, 30], mais selon Légifrance, cette loi du 28 avril 2025 porte sur la sûreté dans les transports [15]. L'état de la loi est à recontrôler le jour de la remise.

\subsection*{Obligations concrètes : NIS2 et RGPD}
\begin{xltabular}{\linewidth}{@{}P{22mm}YY@{}}\toprule
\thd{Thème} & \thd{NIS2 (directive, art.~20, 21, 23, 34)} & \thd{RGPD (art.~32 à 34, 83)}\\\midrule\endfirsthead
\toprule \thd{Thème} & \thd{NIS2 (directive, art.~20, 21, 23, 34)} & \thd{RGPD (art.~32 à 34, 83)}\\\midrule\endhead
\bottomrule\endfoot
Gouvernance & La direction approuve les mesures de gestion des risques et supervise leur mise en œuvre. Elle peut être tenue responsable. Ses membres suivent une formation (art.~20) & Le responsable du traitement, ici MétalTech, prouve sa conformité : registre des traitements, contrats avec les sous-traitants. Le DPO (externe, p.~27) conseille et contrôle, mais l'entreprise reste responsable\\
Mesures minimales & Mesures adaptées au risque et proportionnées, selon une approche \og tous risques\fg{}. Au minimum : analyse des risques, gestion des incidents, continuité et sauvegardes, sécurité de la chaîne d'approvisionnement, sécurité des achats et de la maintenance, évaluation de l'efficacité, hygiène informatique et formation, cryptographie, contrôle d'accès et gestion des actifs, MFA s'il y a lieu (art.~21) & Mesures adaptées au risque : pseudonymisation et chiffrement, confidentialité, intégrité, disponibilité et résilience, capacité à rétablir les données, tests réguliers (art.~32)\\
Notification & Incident important : perturbation grave des services, pertes financières ou dommages considérables à autrui, réels ou possibles. Alerte sous 24~h, notification sous 72~h dès connaissance, sans retard injustifié. Rapport intermédiaire sur demande, final un mois après notification (ou avancement puis final un mois après traitement). (art.~23) & Violation de données personnelles : notification à la CNIL dans les meilleurs délais et, si possible, dans les 72~h après la prise de connaissance, sauf si la violation est peu susceptible de créer un risque. Tout retard doit être justifié. Information des personnes dans les meilleurs délais si le risque est élevé, sauf exceptions (art.~33--34)\\
Preuves & Politiques, analyse des risques, tests, journal des incidents & Registre de toutes les violations, même non notifiées (art.~33.5)\\
Sanctions & Pour une entité importante qui viole l'art.~21 ou 23, chaque État doit prévoir un plafond d'amende au moins égal au plus élevé de ces deux montants : 7~M€ ou 1,4~\% du CA mondial de l'exercice précédent (art.~34.5). Ce n'est ni une amende automatique ni une règle déjà applicable en France & Deux niveaux (art.~83) : jusqu'à 10~M€ ou 2~\% du CA mondial, notamment pour la sécurité et la notification (art.~32 à 34). Jusqu'à 20~M€ ou 4~\% pour d'autres manquements, comme les principes ou les droits des personnes. Le montant le plus élevé est retenu\\
\end{xltabular}
\figsource{Articles lus sur EUR-Lex [10, 11] et page de la CNIL [12], consultés les 24 et 25/11/2026. L'autorité, le portail et les modalités françaises de notification NIS2 dépendront de la loi de transposition. Le PCA et le PAR prévoient déjà l'ANSSI et son centre de réponse aux incidents, le CSIRT [8, §7.3] [9, §8.1].}
\begin{figure}[htp]
\centering
\begin{tikzpicture}[x=1cm,y=1cm]
\draw[thick,graytext,-{Latex[length=2mm]}] (0,0)--(13.6,0);
\foreach \x/\t in {0/{Connaissance\\de l'incident},2.6/{24 h},5.6/{72 h},11.4/{1 mois après\\la notification}}{
  \draw[graytext] (\x,0.12)--(\x,-0.12); \node[font=\scriptsize,align=center,anchor=north] at (\x,-0.18) {\t};}
\node[font=\scriptsize\bfseries,text=cA,anchor=east] at (-0.15,1.25) {NIS2};
\node[font=\scriptsize\bfseries,text=cC,anchor=east] at (-0.15,-1.75) {RGPD};
\foreach \x/\t in {2.6/{alerte précoce\\(CSIRT ou autorité)},5.6/{notification\\d'incident},8.5/{rapport intermédiaire\\(sur demande)},11.4/{rapport final\\(ou d'avancement\\si l'incident dure)}}{
  \node[circle,fill=cA,inner sep=2pt] at (\x,0.45) {};
  \node[box,draw=cA,fill=white,font=\scriptsize,text width=25mm] at (\x,1.35) {\t};}
\node[circle,fill=cC,inner sep=2pt] at (5.6,-0.95) {};
\node[box,draw=cC,fill=white,font=\scriptsize,text width=30mm] at (5.6,-1.8) {notification CNIL si risque\\(si possible sous 72~h)};
\node[box,draw=cC,fill=white,font=\scriptsize,text width=42mm] at (11.0,-1.8) {personnes concernées : dans les meilleurs délais si risque élevé. Registre interne pour toute violation};
\end{tikzpicture}
\caption{Deux procédures distinctes de notification pour un même incident.}
\label{fig:notif}
\figsource{D'après la directive NIS2, art.~23 [10] et le RGPD, art.~33--34 [11, 12]. Axe non proportionnel. Un rançongiciel n'est pas automatiquement une violation à notifier à la CNIL : il faut analyser les données personnelles touchées et le risque pour les personnes.}
\end{figure}

\subsection*{Ce que cela demande à MétalTech}
Les mesures de l'article~21 sont déjà traitées dans ce rapport : analyse des risques et efficacité des mesures (9.2.2), incidents, continuité et sauvegardes (9.2.4), chaîne d'approvisionnement avec l'accès des prestataires par le bastion et des clauses de sécurité dans les contrats (9.2.1, 9.2.6), vulnérabilités (plan sur 18 mois, 9.2.2), cryptographie, accès et actifs (9.2.3), MFA et communications de secours (9.2.6, 9.2.4). Il reste à ajouter la formation des dirigeants (art.~20), des exercices d'hameçonnage avec la cible du cours de moins de 5~\% de clics [2, diap.~25], et une procédure commune DPO--RSSI de notification. Le registre des violations est tenu par l'entreprise, avec l'aide du DPO.

\subsection*{Actions prioritaires vers la certification ISO 27001}
ISO/IEC 27001{:}2022 fixe les exigences d'un SMSI (système de management de la sécurité de l'information). ISO 27002 détaille 93 mesures en 4 thèmes [5, diap.~20]. La certification ne prouve pas seule la conformité au RGPD ou à NIS2, mais elle fournit une grande partie des preuves.

Le sujet vise une certification au T4 2026. Le cas cite un audit à cette date (p.~27), mais sa feuille de route vise 2028 (p.~35). Un auditeur attend un SMSI qui fonctionne déjà, avec au moins un audit interne et une revue de direction. D'où deux scénarios :
\begin{itemize}
\item \textbf{Scénario A, SMSI déjà mature} : si les preuves des étapes 1 à 7 du tableau existent fin 2026, l'audit du T4 2026 reste possible.
\item \textbf{Scénario B, maturité insuffisante} : fin 2026 sert à un pré-audit (diagnostic des écarts par un regard extérieur), puis la certification suit le calendrier ci-dessous, au plus tard à l'horizon 2028.
\end{itemize}

\begin{xltabular}{\linewidth}{@{}C{12.5mm}YP{30mm}P{27mm}@{}}\toprule
\thd{Étape} & \thd{Action et preuve attendue} & \thd{Pilote} & \thd{Échéance (scénario B)}\\\midrule\endfirsthead
\toprule \thd{Étape} & \thd{Action et preuve attendue} & \thd{Pilote} & \thd{Échéance (scénario B)}\\\midrule\endhead
\bottomrule\endfoot
1 & Faire un état des lieux et définir le périmètre du SMSI (site de Louhans, ERP, MES, PLM, réseau, cloud Azure). Preuve : engagement écrit de la direction & PDG, RSSI & Déc. 2026\\
2 & Rédiger la politique de sécurité. Preuve : politique signée & RSSI & Déc. 2026\\
3 & Analyser les risques (méthode de 9.2.2) et écrire le plan de traitement. Preuve : registre des risques & RSSI, responsable SI & Janv. et févr. 2027\\
4 & Rédiger la déclaration d'applicabilité (DdA) : mesures retenues ou exclues, avec justification & RSSI & Févr. 2027\\
5 & Mettre en œuvre les mesures prioritaires et garder les preuves : revues d'accès, tests de restauration, journaux, formations & SI, métiers & Mars à juin 2027\\
6 & Audit interne indépendant, puis revue de direction. Preuves : rapport d'audit, compte rendu de revue & Auditeur interne ou cabinet, direction & Juin et juil. 2027\\
7 & Corriger les non-conformités et choisir un organisme certificateur & RSSI, direction & Sept. 2027\\
8 & Audit de certification par un organisme indépendant & Direction, RSSI & T4 2027, sous condition\\
\end{xltabular}
Nous ne promettons pas la certification : elle dépend des preuves réunies. Le suivi se fait en cycle PDCA : planifier, réaliser, vérifier, améliorer.
\source{[1], p.~27, 34--35 · [2], diap.~25, 28, 30 · [5], diap.~20--21 · [8], §7.3 · [9], §8 · [10] à [15].}
% ---------------------------------------------------------------------
\question{9.2.6}{Sécurisation des accès distants et VPN}{L'entreprise utilise un VPN SSL Fortinet et dispose d'un VLAN invités, dans un contexte hybride avec le cloud Azure. Décrivez les risques associés à cette architecture d'accès distant (phishing, vol de credentials, mouvement latéral) et proposez une stratégie d'authentification renforcée (MFA, Zero Trust, filtrage par rôle) pour sécuriser les connexions des collaborateurs et prestataires externes.}

\subsection*{Situation et risques}
MétalTech utilise un VPN SSL Fortinet pour 3 utilisateurs réguliers [1, p.~26]. Un VPN (réseau privé virtuel) crée un tunnel chiffré entre un poste distant et le réseau de l'entreprise. L'entreprise a aussi un VLAN invités et un annuaire hybride : Active Directory sur site et Azure AD, aujourd'hui Microsoft Entra ID, dans le cloud (p.~26, 28). Le MFA n'est cité que pour les \og accès sensibles\fg{} (p.~26). Les prestataires SAP et MES assurent le support de niveau 2 (p.~31) et ont besoin d'accès à distance. Tous ces accès relèvent de l'IAM, la gestion des identités et des accès [6, diap.~2--4].

\begin{xltabular}{\linewidth}{@{}P{32mm}YP{40mm}@{}}\toprule
\thd{Risque} & \thd{Scénario chez MétalTech} & \thd{Effet}\\\midrule\endfirsthead
\toprule \thd{Risque} & \thd{Scénario chez MétalTech} & \thd{Effet}\\\midrule\endhead
\bottomrule\endfoot
Hameçonnage & Un faux portail VPN ou Microsoft vole l'identifiant d'un commercial ou d'un prestataire & Accès avec un compte valide\\
Vol d'identifiants & Mot de passe réutilisé ou volé, sans MFA sur le VPN & Connexion invisible au départ\\
Mouvement latéral & Si le VPN ouvre un segment entier, un compte volé va plus loin, selon le filtrage et les droits [4, diap.~56] & Propagation vers les serveurs, puis l'atelier\\
Faille du VPN & Boîtier non mis à jour et exposé à Internet (9.2.2) & Authentification contournée\\
Comptes de prestataires permanents & Comptes partagés ou jamais supprimés & Accès après la mission\\
Annuaire hybride & Compromission de l'annuaire local propagée au cloud, ou l'inverse & Perte de contrôle des identités\\
VLAN invités & Mauvaise isolation entre invités ou vers le réseau interne & Entrée d'un visiteur malveillant\\
\end{xltabular}
Un VPN protège le tunnel, mais il ne prouve ni que l'utilisateur est légitime ni que son poste est sain : \og le VPN ne sécurise pas un poste compromis\fg{} [4, diap.~50].

\subsection*{Stratégie proposée : vérifier chaque accès}
\begin{figure}[htp]
\centering
\begin{tikzpicture}[
  u/.style={box,fill=white,text width=30mm,minimum height=10mm},
  c/.style={box,fill=pale,text width=40mm,minimum height=12mm},
  r/.style={box,fill=white,text width=34mm,minimum height=10mm},
]
\node[u] (col) at (0,1.2) {Collaborateurs nomades\\(VPN : 3 réguliers)};
\node[u] (pre) at (0,-0.4) {Prestataires SAP, MES,\\maintenance};
\node[u,fill=rModere!60] (inv) at (0,-2.6) {Invités sur site\\(Wi-Fi)};
\node[c] (id) at (4.1,0.4) {\textbf{1. Identité + MFA}\\AD / Entra ID, MFA\\résistant à l'hameçonnage};
\node[c] (pos) at (4.1,-1.35) {\textbf{2. État du poste}\\EDR actif, correctifs, chiffrement};
\node[c] (gw) at (8.55,-0.45) {\textbf{3. Passerelle}\\VPN SSL à court terme,\\ZTNA (par application) à terme};
\node[c] (pol) at (8.55,-2.55) {\textbf{4. Filtrage par rôle}\\groupes, horaires, durée,\\refus par défaut};
\node[r] (app) at (12.75,0.4) {Applications autorisées\\ERP, CRM, fichiers};
\node[r] (bas) at (12.75,-1.35) {Bastion (PAM)\\sessions enregistrées};
\node[r] (ot) at (12.75,-3.05) {DMZ industrielle $\rightarrow$\\équipement OT sur demande};
\node[r,fill=lightgray] (net) at (4.1,-3.35) {Internet uniquement};
\draw[arr] (col.east)--(id.west); \draw[arr] (pre.east)--(id.west);
\draw[arr] (id)--(pos); \draw[arr] (pos.east)--(gw.west); \draw[arr] (gw)--(pol);
\draw[arr] (pol.east)--(app.west); \draw[arr] (pol.east)--(bas.west); \draw[arr] (bas)--(ot);
\draw[arr] (inv.east)--(net.west);
\draw[arrX] (net.east)--node[below=2pt,font=\tiny\bfseries,text=alert,pos=0.45]{interdit} ($(pol.west)+(0,-0.45)$);
\node[box,draw=cD,fill=cD!8,text width=44mm,font=\scriptsize] (log) at (8.55,1.55) {Journaux centralisés : connexions, refus, sessions (SOC externalisé en projet)};
\draw[cD,dashed,thick] (log.east) -| (app.north);
\draw[cD,dashed,thick] (log.west) -| (id.north);
\end{tikzpicture}
\caption{Accès distant cible : chaque connexion est vérifiée avant d'atteindre une ressource.}
\label{fig:acces}
\figsource{Proposition du groupe, d'après [4, diap.~50, 54--57] et [6, diap.~30--37]. ZTNA : accès réseau Zero Trust, application par application.}
\end{figure}

\textbf{VPN ou Zero Trust ?} Avec un VPN classique, l'utilisateur connecté reçoit souvent un accès large, selon le filtrage et les droits configurés. Le modèle Zero Trust ne fait jamais confiance par défaut : il vérifie l'identité, le poste et le contexte à chaque accès, et limite les droits aux ressources nécessaires [4, diap.~55--56] [2, diap.~13]. Le ZTNA applique ce principe application par application. Pour une PME, nous proposons une transition : renforcer le VPN existant, puis tester le ZTNA sur une application.

\textbf{Prérequis à vérifier.} L'accès conditionnel d'Entra ID demande des licences adaptées, et le contrôle de l'état des postes un outil de gestion des appareils (MDM). L'intégration du VPN Fortinet avec Entra ID doit être testée, par exemple avec SAML, un protocole qui permet à un service de vérifier une identité auprès d'un autre. Le cas ne dit pas si ces licences et outils existent : nous ne les considérons pas comme acquis.

\subsection*{Filtrage par rôle}
\begin{xltabular}{\linewidth}{@{}P{34mm}YP{42mm}@{}}\toprule
\thd{Profil} & \thd{Ressources autorisées} & \thd{Conditions}\\\midrule\endfirsthead
\toprule \thd{Profil} & \thd{Ressources autorisées} & \thd{Conditions}\\\midrule\endhead
\bottomrule\endfoot
Salarié nomade (commercial) & CRM Dynamics 365 (cloud), client ERP, partage commercial & MFA, poste de l'entreprise conforme\\
Direction & ERP en lecture, Power BI, messagerie & MFA, aucun accès d'administration\\
Administrateur SI & Consoles d'administration par le bastion uniquement & Compte d'administration séparé, MFA, sessions enregistrées\\
Prestataire SAP & Serveur ERP par le bastion & Compte nominatif, ouvert sur demande, expiration automatique\\
Prestataire MES & Équipements de la DMZ ou de l'OT par le bastion & Accord du chef d'atelier, créneau défini, session enregistrée\\
Invité & Internet uniquement & Isolation entre appareils Wi-Fi, aucune route vers IT ou OT\\
\end{xltabular}

\subsection*{Mesures et vérifications}
Le PAM (gestion des accès à privilèges) protège les comptes qui ont des droits d'administration [6, diap.~36].

\begin{xltabular}{\linewidth}{@{}P{32mm}YP{40mm}@{}}\toprule
\thd{Mesure} & \thd{Mise en place} & \thd{Vérification}\\\midrule\endfirsthead
\toprule \thd{Mesure} & \thd{Mise en place} & \thd{Vérification}\\\midrule\endhead
\bottomrule\endfoot
MFA pour tous les accès distants & Clés FIDO2 ou clés d'accès (passkeys), qui résistent à l'hameçonnage, si le VPN et les applications sont compatibles. Sinon, application d'authentification. Pas de SMS seul. Récupération sécurisée [6, diap.~6--7] & 100~\% des comptes distants avec MFA\\
Accès conditionnel & Règles Entra ID : poste conforme, pays autorisés, blocage des connexions à risque, sous réserve des licences & Rapport des connexions refusées\\
Moindre privilège & Groupes par rôle, refus par défaut, revue trimestrielle [6, diap.~30, 32] & Compte rendu de revue\\
Comptes de prestataires & Nominatifs, temporaires, demandés par un responsable métier, supprimés en fin de mission [6, diap.~35] & Aucun compte sans date de fin\\
PAM et bastion & Coffre de mots de passe, accès accordé au moment du besoin, sessions enregistrées [6, diap.~36] & Sessions consultables\\
VPN à jour & Firmware FortiGate corrigé, portail VPN exposé seulement si nécessaire & Version vérifiée chaque mois\\
Révocation & Désactivation le jour du départ ou de la fin de mission, sessions coupées & Moins de 24~h\\
\end{xltabular}

\subsection*{Plan et tests}
\textbf{M0--M1} : MFA sur le VPN pour tous, nettoyage des comptes distants et des comptes de prestataires (SI). \textbf{M1--M3} : accès conditionnel si les licences le permettent, bastion pour les prestataires, journaux centralisés (SI, RSSI). \textbf{M3--M6} : pilote ZTNA sur une application, par exemple l'ERP pour les commerciaux, puis décision (SI, direction).

Les tests vérifient, pour chaque profil, les accès autorisés et les accès refusés :
\begin{itemize}
\item un salarié atteint ses applications, mais pas les consoles d'administration ; un administrateur passe obligatoirement par le bastion ;
\item un prestataire n'atteint que son serveur, et son compte expiré est refusé ;
\item un poste non conforme est bloqué, si le contrôle de l'état des postes est disponible ; un invité ne voit aucune adresse interne.
\end{itemize}
Chaque refus doit apparaître dans les journaux.
\source{[1], p.~26--28, 31 · [2], diap.~13 · [4], diap.~50--51, 54--57 · [6], diap.~2--7, 30--37.}
% =====================================================================
\section*{Conclusion et synthèse des priorités}
\addcontentsline{toc}{section}{Conclusion et synthèse des priorités}
Les processus de MétalTech dépendent fortement de son SI : l'ERP, le MES, le PLM et les données qualité servent à fabriquer, à tracer et à livrer. La sécurité concerne donc aussi l'atelier, les achats, la qualité, le commercial et la direction. Notre plan est progressif : réduire d'abord ce qui peut arrêter la production ou exposer les plans des clients, puis améliorer la performance.

\begin{xltabular}{\linewidth}{@{}C{15mm}YP{34mm}P{26mm}@{}}\toprule
\thd{Priorité} & \thd{Action} & \thd{Pilote} & \thd{Échéance}\\\midrule\endfirsthead
\toprule \thd{Priorité} & \thd{Action} & \thd{Pilote} & \thd{Échéance}\\\midrule\endhead
\bottomrule\endfoot
P1 & Copie de sauvegarde hors ligne, test de restauration, fiche réflexe rançongiciel & Responsable SI, RSSI & 0--1 mois\\
P1 & MFA sur tous les accès distants, nettoyage des comptes de prestataires & Responsable SI & 0--1 mois\\
P1 & Inventaire des postes Windows 10 et des firmwares, correctifs critiques & Responsable SI & 0--1 mois\\
P1 & Segmentation IT/OT : pare-feu OT, DMZ industrielle, VLAN de l'atelier & SI, production & 1--8 mois\\
P1 & Trajectoire de l'ERP décidée, support assuré avant fin 2027 & DSI, DAF & Décision à M3\\
P1 & Classification et droits d'accès des plans A320 & Bureau d'études, RSSI & 0--3 mois\\
P1 & Procédure de notification NIS2 et RGPD, préparation de l'ISO 27001 & Direction, RSSI, DPO & 0--3 mois\\
P2 & Pilote SMED et suivi des devis dans le CRM & Production, commercial & 1--6 mois\\
P2 & Deuxième source pour le titane, MRP fiabilisé & Achats, qualité & 3--12 mois\\
\end{xltabular}

\subsection*{Budget par priorité}
Le cas demande une cohérence entre investissement, financement et rentabilité [1, p.~42]. Nous distinguons coûts connus, estimations et devis à obtenir, sans inventer ni gain ni financement. Les montants annuels et pluriannuels ne sont pas additionnés.

\begin{xltabular}{\linewidth}{@{}C{15mm}YP{38mm}P{40mm}@{}}\toprule
\thd{Priorité} & \thd{Poste} & \thd{Montant} & \thd{Statut}\\\midrule\endfirsthead
\toprule \thd{Priorité} & \thd{Poste} & \thd{Montant} & \thd{Statut}\\\midrule\endhead
\bottomrule\endfoot
P1 & SOC externalisé & environ 15~k€ par an & Coût donné par le cas, projet [1, p.~33]\\
P1 & Renouvellement du matériel (postes, réseau) & 25~k€ par an & Budget du cas [1, p.~33]\\
P1 & ESU Windows 10, pour les exceptions seulement & 61~USD par poste la 1\textsuperscript{re} année, prix doublé chaque année (122 puis 244~USD) & Microsoft [16]. Achat par année, 3 ans au plus. Programme cumulatif : une entrée tardive fait payer les années précédentes\\
P1 & Remplacement du pare-feu et des commutateurs & à chiffrer & Devis d'intégrateur. Renouvellement prévu en 2026 [1, p.~33]\\
P1 & Pare-feu OT, bastion, clés FIDO2 & à chiffrer & Devis à obtenir\\
P1 & Assistance à la réponse à incident & à vérifier & Contrat AXA, sinon devis\\
P1 & Projet ERP & 350 à 500~k€ sur 3 ans & Estimation du cas [1, p.~33]\\
P1 & Audit et certification ISO 27001 & à chiffrer & Devis d'organismes certificateurs\\
P2 & Pilote SMED & à estimer & Temps interne, chariot, formation\\
P2 & Qualification d'une deuxième source de titane & à estimer & Audit qualité, échantillons\\
P2 & DLP, accès conditionnel, ZTNA & à vérifier & Selon les licences existantes\\
\end{xltabular}

\textbf{Financement.} Les 28~k€ annuels de sécurité sont inclus dans les 462~k€ du budget SI, et la part encore libre n'est pas connue : il faut vérifier les dépenses engagées (licences, audits, assurance) avant d'en ajouter. Toute nouvelle dépense est votée par la direction. Le cas annonce pour 2025 une capacité d'autofinancement de 1,65~M€ et une trésorerie nette de 2,85~M€ [1, p.~11], mais plusieurs tableaux financiers sont incohérents (annexe A) : ils sont à rapprocher avec la direction financière avant tout engagement.

\textbf{Rentabilité.} Nous ne promettons pas de retour sur investissement chiffré. Ces dépenses réduisent surtout le risque d'un arrêt long, dont le coût dépend de la marge perdue, des pénalités et du rattrapage. La cyberassurance (2~M€) ne remplace pas la prévention.

Chaque action a un responsable, une date, un indicateur et une preuve : rapport de restauration, test de segmentation, revue des droits, tableau de bord du TRS, suivi des devis. Les risques restants sont présentés à la direction pour décision. Une amélioration progressive et mesurée est plus réaliste qu'un changement complet du SI en une fois, avec une priorité : protéger la production, la qualité des pièces et la confiance des clients.
\source{[1], p.~11, 27, 32--35, 42 · [2] · [5] · [7] à [9] · [16].}

% =====================================================================
\section*{Annexe A -- Écarts relevés dans le cas et hypothèses retenues}
\addcontentsline{toc}{section}{Annexe A -- Écarts du cas et hypothèses retenues}
Le cas recommande de faire des hypothèses explicites quand des informations manquent [1, p.~42]. Le tableau ci-dessous liste les contradictions relevées dans les documents originaux et le choix fait dans ce rapport. Les pages renvoient au cas [1], sauf mention contraire.

{\small\begin{xltabular}{\linewidth}{@{}P{26mm}YYP{40mm}@{}}\toprule
\thd{Sujet} & \thd{Version A} & \thd{Version B} & \thd{Hypothèse retenue}\\\midrule\endfirsthead
\toprule\thd{Sujet} & \thd{Version A} & \thd{Version B} & \thd{Hypothèse retenue}\\\midrule\endhead
\bottomrule\endfoot
ERP & SAP Business One déployé en 2019 (p.~19). Version actuelle : 10.0. Date de mise à niveau non précisée. & SAP S/4HANA (p.~6, 7, 12, 40). Nouvel ERP en 2020 (p.~4) & SAP Business One 10.0. S/4HANA est la cible de migration (p.~33)\\
CRM & Dynamics 365 Sales (p.~17, 22) & Salesforce (p.~7, 35, 40) & Dynamics 365\\
CAO et PLM & Windchill et Creo (p.~17, 21) & Siemens NX (p.~6, 40) & Windchill pour le PLM\\
MES & Logiciel en ligne, SaaS (p.~20, 30 et [9, §2.1]) & Sur site (p.~28) & Hébergement à confirmer. Connecteurs locaux dans tous les cas\\
CA 2025 & 18,4~M€ (p.~4, 11, 36) & 13,2~M€ (p.~32) & 18,4~M€. Budget SI = 2,51~\% du CA\\
Bilan 2025 & Total de l'actif : 13\,300~k€ (p.~11) & Somme des lignes du passif : 11\,300~k€ (p.~11) & Bilan déséquilibré. Pas de calcul financier fin sans vérification\\
Marge brute & 7,36~M€, soit 40~\% (p.~11, 36) & CA : 18,4~M€. Achats : 7,2~M€ (p.~11). Budget 2026 : marge brute de 57,5~\% (p.~36). & 40~\% : valeur annoncée. Écart à expliquer selon les stocks et les coûts retenus.\\
Compte de résultat & Total des charges : 18\,050~k€ (p.~12) & Total des produits : 18\,605~k€ (p.~12) & Écart de 555~k€ non expliqué\\
Équipe SI & DSI, 3 administrateurs et RSSI externe à mi-temps (p.~13) & Responsable SI, administrateur SI, poste vacant (p.~30). \og 3 personnes\fg{} (p.~31) & Scénario du sujet : 3 personnes, sans SOC\\
Postes de travail & 45 postes Windows 11 Pro et 15 portables (p.~28, 30) & Windows 10 (p.~33 et énoncé 9.2.2) & Inventaire préalable. Hypothèse du sujet pour 9.2.2\\
Machines de mesure & 1 machine Zeiss listée (p.~6) & 3 CMM connectées (p.~20, 28) & 3 CMM, selon les données d'intégration\\
ISO 27001 & Audit prévu au T4 2026 (p.~27, 34) & Certification en 2028 (p.~35) & T4 2026 est un objectif, pas un acquis\\
NIS2 & Évaluation en cours (p.~27). \og Probable\fg{} (p.~34) & Entité importante (énoncé 9.2.5) & Hypothèse du sujet\\
NADCAP & Prochain audit en décembre 2026 (p.~16) & Accréditation en cours de renouvellement (p.~25) & Renouvellement à suivre\\
Écarts sans effet sur nos réponses & WMS : SAP WM et Easy WMS (p.~10). ISO 9001 acquise en 2010 (p.~4). Effectif : 127 salariés (p.~4) & WMS \og (sage)\fg{} (p.~17), Sage X3 [9, §2.1]. ISO 9001 depuis 2012 (p.~4, 16). Organigramme : 128 (p.~13--14) & WMS non tranché. 127 salariés\\
Devis d'exemple & Lot de 500 pièces, sans quantité par ligne (p.~8) & Total HT égal à la somme des 3 prix affichés (p.~8) & Ambiguïté à clarifier (9.1.3)\\
Support de SAP B1 10.0 & Fin 2027 (p.~33, énoncé 9.2.2) & 31/12/2028 selon la page officielle de SAP, date \og susceptible de changer\fg{} [17] & Échéance du sujet. Date réelle à confirmer\\
Transposition de NIS2 & \og Loi n°~2025-379\fg{} [2, diap.~28, 30] & Loi sur la sûreté dans les transports [15]. Projet de loi adopté en première lecture par l'Assemblée le 17/11/2026, non promulgué au 25/11/2026 [14] & Transposition en cours\\
\end{xltabular}}

% =====================================================================
\section*{Sources et références}
\addcontentsline{toc}{section}{Sources et références}
\small
\subsection*{Supports pédagogiques fournis par l'enseignant}
\begin{enumerate}[label={[\arabic*]},leftmargin=*,itemsep=2pt]
\item Disson, E. (2026). \textit{Cas MétalTech Solutions SAS -- Formation M2 SIA / Audit et Cybersécurité des SI}. Version 1, juin 2026, 42 pages. Source principale. Questions p.~38--39, consignes p.~39, conseils p.~42.
\item Disson, E. (2026). \textit{Fondamentaux en cybersécurité}. Support du 15 septembre 2026, 33 diapositives.
\item Disson, E. (2026). \textit{Le chiffrement : fondamentaux, applications et réglementations}. Support du 22 septembre 2026, 54 diapositives.
\item Disson, E. (2026). \textit{Comprendre, concevoir et défendre les réseaux}. Présentation PowerPoint, 61 diapositives.
\item Disson, E. (2026). \textit{Fondamentaux de l'audit et de la gouvernance des systèmes d'information}. Support du 19 octobre 2026, 36 diapositives.
\item Disson, E. (2026). \textit{Identity \& Access Management (IAM)}. Support du 6 octobre 2026, 39 diapositives.
\item Disson, E. (2026). \textit{Gestion de crise en cybersécurité}. Support du 19 octobre 2026, 39 diapositives.
\item \textit{Plan de continuité d'activité (PCA) -- MétalTech Solutions SAS}, réf. PCA-MTL-001, version 1.0, 21 septembre 2026, 9 pages. Cadre de structuration ISO 22301{:}2019.
\item \textit{Plan d'assurance des reprises (PAR) -- MétalTech Solutions SAS}, réf. PAR-MTL-001, version 1.0, 21 septembre 2026, 11 pages. Fichier intitulé PRA. Cadre ISO/IEC 27031{:}2011.
\end{enumerate}
\subsection*{Sources officielles consultées les 24 et 25 novembre 2026}
\begin{enumerate}[label={[\arabic*]},start=10,leftmargin=*,itemsep=2pt]
\item Union européenne. Directive (UE) 2022/2555 du 14 décembre 2022 (NIS2), art.~20, 21, 23, 34 et annexe II. EUR-Lex : \url{https://eur-lex.europa.eu/legal-content/FR/ALL/?uri=CELEX:32022L2555}
\item Union européenne. Règlement (UE) 2016/679 (RGPD), art.~32 à 34 et 83. EUR-Lex : \url{https://eur-lex.europa.eu/eli/reg/2016/679/oj?locale=fr}
\item CNIL. \textit{Violations de données personnelles : les règles à suivre}, page publiée le 19 juin 2018 : \url{https://www.cnil.fr/fr/violations-de-donnees-personnelles-les-regles-suivre}
\item ANSSI -- MesServicesCyber. \textit{Directive NIS 2} : transposition en cours, pré-enregistrement des entités, référentiel ReCyF : \url{https://messervices.cyber.gouv.fr/nis2}
\item Assemblée nationale. \textit{Adoption du projet de loi relatif à la résilience des infrastructures critiques et au renforcement de la cybersécurité}, 17 novembre 2026 : \url{https://www.assemblee-nationale.fr/dyn/actualites-accueil-hub/adoption-du-projet-de-loi-relatif-a-la-resilience-des-infrastructures-critiques-et-au-renforcement-de-la-cybersecurite}. Dossier législatif sur Vie publique, modifié le 18/11/2026 : \url{https://www.vie-publique.fr/dossierlegislatif/JORFDOLE000050349138}
\item Légifrance. Loi n°~2025-379 du 28 avril 2025 relative au renforcement de la sûreté dans les transports : \url{https://www.legifrance.gouv.fr/loda/id/JORFTEXT000051528207}
\item Microsoft Learn. \textit{Programme de mises à jour de sécurité étendues (ESU) pour Windows 10}, mis à jour le 12/11/2026 : \url{https://learn.microsoft.com/fr-fr/windows/whats-new/extended-security-updates}
\item SAP Support. \textit{Maintenance \& Release Family -- SAP Business One} : \url{https://support.sap.com/en/offerings-programs/support-small-medium-enterprises/business-one/maintenance.html}
\item Microsoft Learn. \textit{Informations sur les versions de Windows 10} (versions 22H2 et LTSC), mis à jour le 19/11/2026 : \url{https://learn.microsoft.com/fr-fr/windows/release-health/release-information}
\item NIST. \textit{SP 800-38D, Recommendation for Block Cipher Modes of Operation: Galois/Counter Mode (GCM) and GMAC}, novembre 2007 (révision annoncée par le NIST en 2024) : \url{https://csrc.nist.gov/pubs/sp/800/38/d/final}
\end{enumerate}
\end{document}
